% !TeX root = detailbericht_de.tex
% !TeX program = pdflatex
\documentclass[11pt,a4paper]{article}
\usepackage[T1]{fontenc}
\usepackage[utf8]{inputenc}
\DeclareUnicodeCharacter{2212}{\ensuremath{-}}
\DeclareUnicodeCharacter{2265}{\ensuremath{\ge}}
\DeclareUnicodeCharacter{00D7}{\ensuremath{\times}}
\usepackage[ngerman]{babel}
\usepackage{lmodern,microtype}
\usepackage[margin=2.1cm,headheight=14pt,footskip=11mm]{geometry}
\usepackage{amsmath,amssymb,graphicx,booktabs,array,longtable,xcolor,caption}
\usepackage[numbers,sort&compress]{natbib}
\usepackage[hidelinks]{hyperref}
\definecolor{blue}{HTML}{155F87}
\definecolor{muted}{HTML}{555555}
\hypersetup{pdftitle={NK/CMV: Deutscher Detail- und Prüfbericht},pdfauthor={SSBI-Projektgruppe},bookmarksnumbered=true}
\graphicspath{{detail_assets/}}
\setlength{\parindent}{0pt}
\setlength{\parskip}{5pt plus 1pt}
\setlength{\textfloatsep}{8pt}
\setlength{\intextsep}{8pt}
\setlength{\tabcolsep}{5pt}
\renewcommand{\arraystretch}{1.16}
\captionsetup{font=small,labelfont=bf,skip=4pt}
\setcounter{tocdepth}{1}
\raggedbottom
\emergencystretch=2em
\newcommand{\code}[1]{\expandafter\nolinkurl\expandafter{\detokenize{#1}}}
\newcommand{\source}[1]{\par{\footnotesize\raggedright\color{muted}\textbf{Nachvollziehen:} #1\par}}
\newcommand{\finding}[1]{\par\smallskip\noindent\textcolor{blue}{\textbf{Prüfbefund.}} #1\par}
\newcommand{\fig}[4]{\begin{figure}[!ht]\centering\includegraphics[width=#2\linewidth]{#1.pdf}\caption{#3}\label{#4}\end{figure}}
\newcommand{\tab}[4]{\begin{table}[!ht]\centering\caption{#2}\label{#3}{\small\input{detail_assets/#1.tex}}#4\end{table}}
\newcommand{\stepbox}[2]{\fbox{\parbox[c]{#1}{\centering\small\vspace{4pt}#2\vspace{4pt}}}}
\newcommand{\AliveCells}{3438750}
\newcommand{\NkCells}{261593}
\newcommand{\HashChecks}{48}
\newcommand{\ProfileChecks}{5180}
\newcommand{\CellcnnFilters}{411}
\newcommand{\CnnEpochMedian}{18,5}
\newcommand{\CnnMaxEpochs}{10}
\newcommand{\NullModels}{15}

\begin{document}
% Seite 1
{\color{blue}\Large SSBI-Gruppenprojekt\par}
\vspace{7mm}
{\LARGE\bfseries NK/CMV-Single-Cell-Analyse\par}
\vspace{3mm}
{\Large Deutscher Detail- und Prüfbericht\par}
\vspace{4mm}
{\large Was wurde gemacht, warum wurde es gemacht, und was lässt sich daraus schließen?\par}
\vspace{7mm}
\textbf{Dokumentationsstand: 6. September 2026.} Grundlage ist der Analysecode des
Commits \code{19432f2} und der vollständige gespeicherte Hauptlauf auf
\code{gated_alive}. Dieser ergänzende Arbeitsbericht erklärt die Analyse;
der englische fünfseitige Bericht bleibt die kompakte Abgabeversion.

\section{Fragestellung und Leseanleitung}\label{sec:frage}
Wir untersuchen, ob Einzelzell-Proteinmessungen aus Blutproben den CMV-Status
bisher ungesehener Spender vorhersagen können. Anschließend fragen wir,
welche Zellen nach den jeweiligen Modellen diese Vorhersage unterstützen.
CMV-Status bedeutet hier die vorgegebene Zuordnung zu CMV-seronegativen bzw.
-seropositiven Spendern. Es gibt \emph{keine unabhängig gemessenen CMV-Labels
für einzelne Zellen} und keine Referenzannotation aller Zelltypen.

Die Daten enthalten 20 Spender, aber Millionen Zellen. Diese große Zellzahl
ersetzt keine große Spenderkohorte. Eine Zelle ist ein Messereignis mit 37
Analysemarkern; ein Spender ist die unabhängige Beobachtung für den Benchmark.
Unser Anspruch ist deshalb eine korrekte spenderweise Auswertung mit
transparenten Grenzen, keine klinische Validierung.

\textbf{Die drei Hauptmethoden beantworten die Frage unterschiedlich.}
CellCNN lernt Zellfilter und Spenderklassifikation gemeinsam. Citrus bildet
zuerst hierarchische Zellgruppen und modelliert deren Häufigkeiten. Die
lineare Single-Cell-SVM lernt eine Zellmarge und verdichtet die höchsten
Margen zu einem Spenderscore. Aufgabe 5 überträgt diese Modelle zurück auf
konkrete Zellen derselben explorativen Karte.

\textbf{Zentrales Ergebnis.} CellCNN und SVM erreichen beide eine mediane
ROC-AUC von 0,875, Citrus 0,625. Daraus folgt keine nachgewiesene Überlegenheit
von CellCNN gegenüber SVM. Die ausgewählten Zellprofile sind methodenspezifisch;
auch ähnliche Markertrends beweisen keine identische biologische Population.

Dieser Bericht erklärt zuerst die gemeinsame Grundlage, dann die drei
Methoden und schließlich die Interpretation. Tabellen und Grafiken sind
Bestandteil der Argumentation: Ihre Bildunterschriften nennen Datenbasis,
Berechnung und Aussagegrenzen. Abweichungen vom Paper werden explizit
benannt; Prüfungen erhalten einen unterscheidbaren Nachweisstatus.
\source{Aufgabenblatt, Projekt 1; \code{CellCNN.pdf}; Notebooks 01--05;
\code{results/tables/task4_metric_summary.csv}.}

\clearpage
% Seite 2
\begingroup\small
\setlength{\parskip}{3pt}
\tableofcontents
\vspace{4mm}
\subsection*{Begriffe, die im ganzen Bericht gelten}
\begin{description}
\item[Zelle / Event:] Eine Zeile der ausgewählten FCS-Markermatrix. Ereignisindex
und Spender-ID erlauben das Wiederfinden der unveränderten Originalmessung.
\item[Spender / Sample:] Alle Zellen einer Person mit einem gemeinsamen CMV-Label.
Die Zuordnung zu Training, Validierung und Test geschieht auf dieser Ebene.
\item[Split und Fold:] Ein äußerer Split teilt die 20 Spender in 14 Trainings-
und sechs Testspender. Ein innerer Fold bestimmt, welche der 14 Spender
innerhalb der Modellauswahl als Validierung dienen.
\item[Score und Wahrscheinlichkeit:] Ein Score ordnet Beobachtungen.
SVM-Margen sind keine Wahrscheinlichkeiten. Auch numerische Softmax- und
Logistikwerte sind ohne zusätzliche Prüfung nicht als kalibrierte Risiken belegt.
\item[OOF:] \emph{Out of fold}, also bewertet von einem Modell, dessen
Trainingsmenge den betreffenden Spender ausschließt. Aufgabe 5 verwendet
strenger ausschließlich dessen \emph{äußere Testmodelle}.
\item[Selektion:] Eine quantitativ definierte Auswahl von Zellen. Sie ist
kein zusätzlich gemessenes Zelltyp-Label und kein Kausalnachweis.
\end{description}
\subsection*{Wie Prüfungen gekennzeichnet werden}
\textbf{Jetzt geprüft} bezeichnet Berechnungen oder Codeabgleiche, die beim
Erstellen dieses Berichts tatsächlich ausgeführt wurden. \textbf{Vorhandener
Prüfnachweis} bezeichnet frühere Rechenprüfungen, deren gespeicherte Artefakte
auf den aktuellen Stand geprüft wurden. \textbf{Offen / nicht verifiziert}
bezeichnet fehlende Belege oder nicht erneut ausgeführte Modellrechnungen.
Ein erfolgreicher Konsistenztest bedeutet nicht automatisch, dass eine
biologische Schlussfolgerung richtig oder das gewählte Verfahren optimal ist.
\endgroup

\clearpage
% Seite 3
\subsection*{Aufgabe 1: Was zeigt das Originalpaper?}
Arvaniti und Claassen entwickeln CellCNN für die Suche nach
phänotypassoziierten Zellsubsets aus ungeordneten Einzelzellmessungen
\cite{cellcnn}. Ein Gesamtlabel, beispielsweise eine Erkrankungsgruppe,
überwacht die Bildung einer passenden Zellrepräsentation. Der Ansatz
verbindet damit \emph{Multiple Instance Learning}: Viele Instanzen (Zellen)
gehören zu einer beschrifteten Menge (Spenderprobe).

Das Paper behandelt mehrere Aufgaben: stimulierte PBMC-Proben, mit
AIDS-freiem Überleben assoziierte T-Zellsubsets, CMV-assoziierte seltene
NK-Subsets und künstliche Leukämie-Spike-ins. Die Resultate dieser
Experimente sind nicht austauschbar. Unser Datensatz betrifft ausschließlich
die CMV-Kohorte; wir reproduzieren weder Überlebensanalysen noch AML-Spike-ins.

\begin{table}[!ht]\centering\small
\caption{Im Paper diskutierte oder verglichene Ansätze und ihre Rolle im Projekt.}
\begin{tabular}{p{3.5cm}p{7.2cm}p{4.5cm}}\toprule
Ansatz & Grundprinzip & Projektstatus\\\midrule
CellCNN & Gelernte Zellfilter, Pooling, überwachte Ausgabe & Hauptmethode\\
Citrus & Hierarchische Cluster, Probenmerkmale, regulärer Prädiktor & Hauptmethode\\
Single-Cell-SVM & Jede Zelle erhält das Label ihrer Probe & Lineare Hauptbaseline mit eigener Spenderaggregation\\
Logistische Regression / Random Forest & Zellbasierte überwachte Klassifikation & Literaturübersicht; nicht zusätzlich trainiert\\
Distanzbasierte Ausreißersuche & Abweichung von Kontrollzellen & Literaturübersicht; nicht implementiert\\
Momente & Erste vier Momente der Marker-Verteilungen als Probenbeschreibung & Optional geplant; nicht als vierte Methode ausgeführt\\
Denoising-Autoencoder & Rekonstruktion gestörter Zellinputs, unüberwachte Repräsentation & Literaturübersicht; nicht implementiert\\
\bottomrule\end{tabular}
\end{table}

\textbf{Warum diese Auswahl?} Die Aufgabenstellung verlangt drei vergleichende
Ansätze. Im Projekt sind CellCNN, Citrus und die lineare Single-Cell-SVM
als Hauptmethoden festgelegt. Der CMV-Vergleich des Papers und der dortige
allgemeine Single-Cell-Baselinevergleich sind nicht exakt dieselbe Auswertung.
Insbesondere unsere Top-1-\%-Aggregation für SVM ist eine Projektanpassung;
das Paper beschreibt nicht diesen vollständigen CMV-SVM-Benchmark.

\textbf{Referenzen sind keine identischen Rezepte.} Die offiziellen
Notebooks \code{NK_cell.ipynb} und \code{NK_cell_ungated.ipynb} zeigen einzelne
Beispielläufe mit eigener Train-/Validierungsaufteilung und Bibliotheksdefaults
\cite{cellcnncode}. Das Paper beschreibt hingegen 100 äußere Wiederholungen
mit innerer dreifacher Auswahl. Bei Widersprüchen haben Aufgabenstellung und
Paper Vorrang. Wir übernehmen keine veraltete Python-/Keras-Umgebung, sondern
implementieren den relevanten Kern in PyTorch neu.
\source{Paper: Methodenabschnitte „Model selection and interpretation“ und
„Baseline methods“; offizielle Notebooks lesend abgerufen, nicht ausgeführt.}

\clearpage
% Seite 4
\section{Daten, Vorverarbeitung und Qualitätskontrolle}\label{sec:daten}
Die bereitgestellten Daten stammen aus der Untersuchung menschlicher
NK-Zelldiversität von Horowitz et al.\cite{horowitz}; das Benchmarkarchiv
wird über Zenodo bereitgestellt\cite{dataset}. Es enthält zwei Gate-Stufen
für dieselben 20 Spender: \code{gated_alive} mit lebenden, von Doubletten
bereinigten PBMCs und \code{gated_NK} nach zusätzlicher NK-Selektion.
AppleDouble-Begleitdateien im macOS-Ordner sind keine Messproben.

Die gemeinsame Labeldatei nennt NK-Dateinamen. Der Loader entfernt den
Gate-Suffix und ordnet anhand der Spender-ID die Alive-Datei zu. Der
Dateiname entscheidet damit nicht versehentlich über das Haupt-Gate.
11 Labels sind 0 (CMV−), neun sind 1 (CMV+). Diese Zuordnung wurde jetzt
gegen sämtliche Split- und Testvorhersagezeilen geprüft.
\tab{spender}{Vollständiges Spenderinventar. Zellzahlen jetzt gegen die FCS-Dateien
geprüft. „Testmodelle“ zählt äußere Testauftritte je Methode über 100 Splits.}{tab:spender}{}

\textbf{Warum gleiche Stichprobengrößen?} Ohne Ausgleich dominieren zellreiche
Proben eine zusammengeführte Zellverteilung. Das Sampling gleicht deshalb
die Zahl der Zellen bzw. Multi-Cell-Inputs pro Spender an. Es gleicht nicht
automatisch die Anzahl der Spender je Klasse in jedem inneren Fold an.

Die 37 Marker stehen in verbindlicher Reihenfolge in \code{NK_markers.csv}.
Technische Kanäle wie Zeit, Ereignislänge, DNA oder Totzellkanal werden nicht
als Analysemarker verwendet. Das Paper nennt 36 Marker; beide offiziellen
Beispielnotebooks enthalten dieselben 37 wie unsere Datei. Diese konkrete
Abweichung der Markerzahl wird nicht durch stilles Entfernen eines Markers
„korrigiert“.
\source{Notebook 01, Abschnitte 2--4; Notebook 04a, Abschnitte 2--4;
\code{untersuchung/inspect_fcs.py}, \code{read_fcs_metadata}.}

\clearpage
% Seite 5
\subsection*{Die anfängliche NK-Exploration}
Notebook 01 untersucht zunächst die kleinere NK-Stufe. Pro Spender werden
1.000 Zellen ohne Zurücklegen gezogen, mit \code{DataFrame.sample} und Seed
\(42+i\), wobei \(i\) die nullbasierte Position in der sortierten Spenderliste
ist. Es entstehen 20.000 Zellen. Dies ist eine explorative Sicht auf bereits
NK-vorselektierte Daten und keine unabhängige Gate-Sensitivitätsanalyse.
\fig{zellzahlen}{.98}{Zellzahlen und Stärke des NK-Gatings. Links logarithmische
Zellzahl je Spender, rechts Anteil der NK-Ereignisse an Alive-Ereignissen.
Jeder schwarze Punkt rechts ist ein Spender; keine Signifikanzprüfung.}{fig:zahlen}
\fig{nk_histogramme}{.93}{Sechs Marker der ursprünglichen NK-Exploration,
1.000 Zellen je Spender, \(\operatorname{arcsinh}(x/5)\). Dichten sind innerhalb
jeder CMV-Gruppe normiert. Die Kurven illustrieren Verteilungen; sie bewerten
keinen Klassifikator und sind kein Test auf unabhängigen Zellen.}{fig:nkhist}
\clearpage
\subsection*{Was die gemeinsame NKG2C-/CD57-Verteilung zeigt}
Die vorherigen Histogramme zeigen jeden Marker einzeln. Zwei erhöhte
Randverteilungen lassen offen, ob dieselben Zellen beide Marker tragen.
Die zweidimensionale Verteilung ergänzt deshalb eine konkrete gemeinsame
Messung. Jeder Punkt der zugrunde liegenden Daten beschreibt die NKG2C-
und CD57-Werte ein und desselben NK-gegateten Ereignisses.

\fig{nk_koexpression}{1}{Gemeinsame NKG2C-/CD57-Verteilung derselben
NK-Stichprobe. Gemeinsame logarithmische Farbskala mit Zellzahlen pro Hexagon.
Hier sind tatsächlich gemeinsame Messungen sichtbar; daraus folgt noch keine
Definition eines validierten CMV-assoziierten Gates.}{fig:nkhex}

Hexagone fassen nahe Messpunkte zusammen; ihre Farbe zeigt Ereigniszahlen,
keine Erkrankungswahrscheinlichkeit. Eine gemeinsame logarithmische Skala
verhindert, dass dieselbe Farbe zwischen den Panels verschiedene Häufigkeiten
bedeutet. Die CMV−-Gruppe enthält 11.000, die CMV+-Gruppe 9.000 Zellen,
weil jede Person gleich viele Zellen beiträgt, aber die Gruppen verschieden
viele Personen umfassen. Die Farben sind daher keine unmittelbar
vergleichbaren gruppennormierten Wahrscheinlichkeitsdichten.

Das Paper berichtet eine NKG2C-/CD57-assoziierte NK-Population. Die Grafik
ist damit eine zielgerichtete Plausibilitätsansicht der Ausgangsdaten.
Sie legt weder Expressionsschwellen fest noch trainiert sie ein Gate.
Aus ihr wurden keine Spender ausgeschlossen und keine Hyperparameter der
Klassifikatoren anhand äußerer Testlabels ausgewählt.

Auch eine sichtbare gemeinsame Expression ist noch kein Nachweis, dass
genau diese Zellen für die spätere Spendervorhersage verantwortlich sind.
Die Grafik verwendet das kleinere NK-Gate und eine andere Zellstichprobe
als Aufgabe 5. Der Hauptvergleich und dessen Interpretation verwenden
Alive-Daten. Die beiden Ansichten dürfen deshalb nicht als wiederholte
Bestätigung derselben ausgewählten Zellpopulation behandelt werden.
\source{Notebook 01, Abschnitte 3, 5 und 6. Die deutsche Neuzeichnung verwendet
identisches Sampling; die Hexagon-Farbskala ist ausdrücklich gemeinsam normiert.}

\clearpage
% Seite 6
\subsection*{Haupt-Gate: kontrolliertes Laden und numerische Transformation}
FlowKit liest die Originaldateien mit \code{ignore_offset_error=True}:
Der bekannte Datenendoffset ist um ein Byte inkonsistent.
\code{source="raw"} liefert untransformierte Werte; die vorhandenen
biologischen Gates bleiben Teil der Datenbasis.

Die feste Transformation lautet
\[
 u=\operatorname{arcsinh}(x/5)
   =\log\!\left(x/5+\sqrt{1+(x/5)^2}\right).
\]
Nahe null ist sie annähernd linear, große Werte werden komprimiert;
negative Werte bleiben erhalten. Kofaktor 5 stammt aus der Referenz,
nicht aus einer Testlabel-Optimierung. Beispiele: \(0\mapsto0\),
\(5\mapsto0{,}8814\), \(-5\mapsto-0{,}8814\) (gerundet).
\fig{qc_marker}{.95}{Alive-QC: 2.000 Zellen je Spender, insgesamt 40.000,
\code{default_rng(12345)}, ohne Zurücklegen. Links Boxplots aller 37
transformierten Marker; ausgeblendete Ausreißersymbole bedeuten keine
Datenentfernung. Rechts Spender-Mediane, je Marker robust skaliert:
\((m_d-\operatorname{median}(m))/(1{,}4826\operatorname{MAD}(m))\).
Die Farbgrenzen bei ±3 betreffen nur die Darstellung.}{fig:qc}

Alle ausgewählten Markerwerte der 40 FCS-Dateien wurden jetzt erneut auf
Endlichkeit geprüft, die Eventzahlen unabhängig aus den Metadaten
gegengelesen. Die Alive-Stufe umfasst 3.438.750 Zellen. Die QC-Heatmap zeigt
mögliche technische Verschiebungen; sie begründet keine automatische
Batchkorrektur. Es werden keine Spender entfernt und keine zusätzliche
Bead-Normalisierung durchgeführt, da Bead-Referenzdaten fehlen.
\source{Notebook 04a, Abschnitte 3--8. Die robuste QC-Skalierung wird nicht an
SVM oder CellCNN weitergegeben. Der Bericht exportiert keine geänderten FCS-Dateien.}

\clearpage
% Seite 7
\section{Aufgabe 2: Dimensionsreduktion}\label{sec:task2}
\subsection*{Eine feste Zellbasis für alle Karten}
Aus \code{gated_alive} werden 500 Zellen je Spender ohne Zurücklegen gewählt,
zusammen 10.000. \code{default_rng(42+i)} erzeugt die Indizes, die danach
sortiert werden. Das ist nicht derselbe Zufallsgenerator wie
\code{DataFrame.sample} in der NK-Vorexploration. Die Zell-ID
\code{gated_alive:Spender:Ereignisindex} hält jede Messung stabil fest.
Alle acht Karten verwenden dieselben IDs und dieselbe Markerreihenfolge.

Nach \(u_{ij}=\operatorname{arcsinh}(x_{ij}/5)\) werden Marker standardisiert:
\[
 z_{ij}=\frac{u_{ij}-\mu_j}{\sigma_j},\quad
 \mu_j=\frac1N\sum_i u_{ij},\qquad
 \sigma_j^2=\frac1N\sum_i(u_{ij}-\mu_j)^2.
\]
\code{StandardScaler} verwendet die Populationsvarianz mit Divisor \(N\).
Hier wird auf allen 20 Spendern gemeinsam gefittet: eine ausdrücklich
vereinbarte Ausnahme für die \emph{explorativen} Aufgaben 2/3. Diese
Skalierung und die Karten sind keine Eingangsdaten der Klassifikatoren.
Für Aufgabe 5 liefern sie nur die Darstellungsebene.

\subsection*{PCA: lineare Varianzachsen}
PCA bestimmt orthogonale Richtungen maximaler Varianz in der 37D-Matrix.
Wir behalten zwei Komponenten, berechnet mit vollständiger SVD. Ohne
Whitening wird auf diese Achsen projiziert. Whitening teilt jede Achse
zusätzlich durch ihre geschätzte Standardabweichung. Die Richtungen bleiben
gleich, die Achsenskalierung und damit euklidische Nachbarschaften können
sich ändern. Die ersten beiden Komponenten erklären hier 16,09 und
10,31\%, zusammen rund 26,39\% der standardisierten Gesamtvarianz.

\subsection*{t-SNE: lokale Nachbarschaftswahrscheinlichkeiten}
t-SNE\cite{tsne} passt eine zweidimensionale Verteilung so an, dass lokale
Nachbarschaften aus dem hochdimensionalen Raum ähnlich repräsentiert werden.
Die Perplexität steuert die effektive Nachbarschaftsgröße. Wir testen 5, 30
und 50, mit PCA-Initialisierung, automatischer Lernrate, 1.000 Iterationen,
Early Exaggeration 12, Barnes-Hut-Näherung mit Winkel 0,5 und Seed 42.
Andere wählbare Parameter sind Distanzmetrik, Initialisierung und Lernrate;
sie werden hier konstant gehalten, nicht zusätzlich durchsucht.

\subsection*{UMAP: Nachbarschaftsgraph und zweidimensionale Anordnung}
UMAP\cite{umap} konstruiert eine gewichtete Nachbarschaftsstruktur und sucht
eine entsprechende Anordnung. Wir vergleichen \((n,d)=(15;0{,}1)\),
\((50;0{,}1)\), \((15;0{,}5)\), mit \(n\) (\code{n_neighbors}) und
\(d\) (\code{min_dist}). Größeres \(n\) berücksichtigt größere Umgebung;
\(d\) beeinflusst die mögliche Verdichtung. Fix sind euklidische Distanz,
spektrale Initialisierung, 500 Epochen, zwei Komponenten und Seed 42.
\source{\code{src/task23_analysis.py}: \code{load_exploratory_data},
\code{projection_configs}, \code{fit_projection}; Notebook 02, Abschnitte 1--2.}

\clearpage
% Seite 8
\subsection*{Alle acht Karten und ihr sichtbarer Parameter-Effekt}
\fig{projektionen}{1}{Alle acht Projektionen der gleichen 10.000 Alive-Zellen.
Farbe: gemessene CD3-Intensität nach ArcSinh, gemeinsame Skala und identische
Punktreihenfolge. Die Karte selbst wird aus allen 37 standardisierten
Markern berechnet. Koordinaten werden unverändert aus
\code{task2_embeddings.csv} geladen; keine neue Anpassung.}{fig:projektionen}

PCA zeigt breite Überlagerungen. Die Nichtlinearität von t-SNE und UMAP
stellt lokale Zellgruppen stärker räumlich getrennt dar. Bei Perplexität 5
zerfällt die t-SNE-Karte stärker in kleine Bereiche; 30 und 50 liefern
weniger stark fragmentierte lokale Strukturen. Bei UMAP ändern sich Form,
Packung und Trennung sichtbar mit Nachbarzahl und Mindestabstand.

Diese Beobachtungen sind Parameterempfindlichkeiten der Darstellung.
\textbf{Eine Insel ist nicht automatisch ein Zelltyp.} Zwischeninselabstände,
Flächengrößen und relative Orientierung sind insbesondere bei t-SNE keine
direkt ablesbaren biologischen Abstände oder Populationshäufigkeiten.
CD3 ist hier eine Interpretationshilfe, keine zur Kartenoptimierung
verwendete Zelltypannotation.

Ein bloßer Koordinatenvergleich wäre ungeeignet: Eine gedrehte oder
verschobene Karte kann dieselben Beziehungen ausdrücken. Deshalb wird im
nächsten Schritt ein orientierungsunabhängiger Vergleich der lokalen
Nachbaridentitäten verwendet. Er beantwortet eine konkrete Frage: Bleiben
für eine Zelle dieselben anderen Zellen in ihrer unmittelbaren Umgebung?

\finding{Die beiden PCA-Varianten unterscheiden sich geringfügig in der
Nachbarschaftsmetrik. „Whitening ändert nichts“ wäre zu stark; korrekt ist
„in diesem Raster nahezu gleicher gemessener Nachbarschaftserhalt“. Die
berichtete Auswahl ungewichtet standardisierter PCA-Achsen folgt den Zahlen.}
\source{Notebook 02, Abschnitt 2; \code{task2_parameters.csv},
\code{task2_provenance.json}, \code{task2_embeddings.csv}.}

\clearpage
% Seite 9
\subsection*{Quantitativer Vergleich und begrenzte Empfehlung}
Für jede Zelle \(i\) bestimmt jede Karte die 15 nächsten \emph{anderen}
Zellen. Mit \(N^A_{15}(i)\) und \(N^B_{15}(i)\) ist
\[
 J_i(A,B)=\frac{|N^A_{15}(i)\cap N^B_{15}(i)|}
 {|N^A_{15}(i)\cup N^B_{15}(i)|}.
\]
Bei zehn gemeinsamen Nachbarn sind es beispielsweise
\(10/(15+15-10)=0{,}5\), nicht \(10/15\). Zuerst wird innerhalb jedes
Spenders gemittelt, anschließend über die 20 Spender. Alle
\(\binom82=28\) Kartenpaare werden bewertet; die Diagonale einer Karte
mit sich selbst beträgt eins.
\fig{nachbarschaften}{.65}{Vollständige Paarmatrix. P/Pw: PCA ohne/mit Whitening;
T5/T30/T50: t-SNE-Perplexität; U15/U50: UMAP-Nachbarzahl bei \(d=0{,}1\);
Ud: UMAP mit \(n=15,d=0{,}5\). Der Jaccard-Index bewertet lokale Geometrie,
keine biologische Richtigkeit.}{fig:jaccard}
\begin{minipage}[t]{.47\linewidth}\small
\begin{tabular}{lr}
\toprule
Einstellung & Jaccard zu 37D \\
\midrule
PCA & 0,01021 \\
PCA mit Whitening & 0,01020 \\
t-SNE: p = 5 & 0,16213 \\
t-SNE: p = 30 & 0,18237 \\
t-SNE: p = 50 & 0,17975 \\
UMAP: n = 15, d = 0,1 & 0,09171 \\
UMAP: n = 50, d = 0,1 & 0,07378 \\
UMAP: n = 15, d = 0,5 & 0,05315 \\
\bottomrule
\end{tabular}

\captionof{table}{Nachbarschaftserhalt gegenüber den standardisierten 37
Markerdimensionen; gleiche Zellen, \(k=15\).}\label{tab:projektion}
\end{minipage}\hfill
\begin{minipage}[t]{.49\linewidth}
Je Verfahren wird die Einstellung mit höchstem 37D-Jaccard gewählt.
Gleichstand innerhalb \(10^{-12}\) bevorzugt die Baseline, danach die
Rasterreihenfolge. Gewählt werden PCA ohne Whitening, t-SNE mit 30 und
UMAP mit \((15;0{,}1)\). Für Aufgabe 5 wird die t-SNE-30-Karte verwendet.

Die kleinen absoluten Werte zeigen Informationsverlust. Die Empfehlung
gilt für dieses Kriterium, diese Stichprobe und dieses Raster, nicht für
alle möglichen Einstellungen oder ein global bestes biologisches Bild.
\end{minipage}
\source{\code{neighbor_indices}, \code{neighbor_jaccard}, \code{recommend_projections};
Notebook 02, Abschnitte 3--5. Exakte euklidische Nachbarn; Eigenidentitäten
werden auch bei Distanzgleichständen entfernt.}

\clearpage
% Seite 10
\section{Aufgabe 3: Clustering und Zelltypen}\label{sec:task3}
Clustering wird in den \textbf{37 standardisierten Markerdimensionen}
durchgeführt, nicht auf einer verzerrenden 2D-Karte. Daten, Sampling und
explorative Skalierung entsprechen Aufgabe 2. Notebook 03 lädt sie
selbstständig und benötigt keinen gemeinsam gehaltenen Notebook-Kernel.

\textbf{K-Means} minimiert \(\sum_i\|z_i-\mu_{c_i}\|^2\). Getestet werden
6, 10 und 14 Cluster, jeweils 20 Initialisierungen, maximal 300 Iterationen
und Seed 42. K-Means bevorzugt kompakte Gruppen im euklidischen Raum.
\textbf{Ward} verbindet schrittweise diejenigen Gruppen, deren Vereinigung
die innerhalb der Gruppen liegende Quadratsumme am wenigsten vergrößert.
Ein gemeinsamer Hierarchiebaum wird bei 6, 10 und 14 Clustern geschnitten.

\textbf{Leiden}\cite{leiden} arbeitet auf einem Graphen: Eine Zelle wird mit
ihren 15 nächsten Nachbarn verbunden. Die Vereinigung beider Richtungen
bildet einen ungewichteten, ungerichteten Graphen ohne Selbstkanten und
Mehrfachkanten. Der resultierende Grad muss nicht 15 sein. Optimiert wird
\code{RBConfigurationVertexPartition}, mit Auflösungen 0,5, 1,0, 1,5,
Seed 42 und Iteration bis zum Stillstand (\code{n_iterations=-1}). Die
Auflösung ist keine fest vorgegebene Clusterzahl.
\tab{clustervergleich}{Alle neun Lösungen. Silhouetten auf einer gemeinsamen
Stichprobe von 100 Zellen je Spender, insgesamt 2.000. Mindestgrößen betreffen
hingegen die vollständigen 10.000 explorativen Zellen.}{tab:cluster}{}

\textbf{Silhouette als Auswahlkriterium.} Für eine Zelle ist \(a_i\) die
mittlere Distanz zu anderen Zellen desselben Clusters, \(b_i\) das kleinste
mittlere Distanzmaß zu einem anderen Cluster. Dann
\[
 s_i=\frac{b_i-a_i}{\max(a_i,b_i)}\in[-1,1].
\]
Positive Werte sprechen für größere Nähe zur eigenen Gruppe. Die Distanzen
werden hier innerhalb der gemeinsamen 2.000-Zell-Stichprobe berechnet;
es ist eine Approximation der Vollstichprobensilhouette, nicht nur ein
Mittel ausgewählter exakt berechneter Vollstichprobenwerte.

Die Auswertungszellen werden mit Seed \(43+i\) gewählt. Falls einem
Verfahren dadurch ein Cluster vollständig fehlen würde, stellt die
Implementierung \emph{alle} Lösungen auf die Gesamtauswertung um. Im
vorliegenden Lauf war dies nicht nötig. Pro Methode gewinnt die höchste
spendergemittelte Silhouette; Gleichstand innerhalb \(10^{-12}\) bevorzugt
weniger Cluster, danach die Rasterreihenfolge.
\source{\code{fit_clusterings}, \code{make_neighbor_graph},
\code{silhouette_evaluation}, \code{select_clusterings}; Notebook 03.}

\clearpage
% Seite 11
\subsection*{Clusterkarten und die zugehörigen Markerprofile}
\fig{clusterkarten}{1}{Ausgewählte Clusterlösungen auf derselben gespeicherten
UMAP-Basiskarte (15 Nachbarn, Mindestabstand 0,1). Die Farben kodieren nur
methodenspezifische IDs; gleiche Farbe bedeutet nicht denselben Zelltyp.
Clusterbildung im 37D-Raum, Projektion nur zur Darstellung.}{fig:clusterkarten}
\fig{clusterprofile}{1}{Markerprofile der ausgewählten Lösungen für zwölf
interpretierbare Marker. Je Spender und Cluster wird der Marker-Median
berechnet, anschließend der Median über vertretene Spender. Danach wird
die gespeicherte explorative z-Skala angewandt. Gemeinsame Farbskala,
keine neue Standardisierung.}{fig:clusterprofile}

Fehlende Spender-Cluster-Kombinationen werden bei Profilen nicht als
Null-Expression imputiert. Für Clusterhäufigkeiten ist dagegen ein fehlender
Cluster in einer Probe korrekt eine Häufigkeit null. Diese beiden Fälle
sind fachlich verschieden. Zusätzlich werden Zellzahl, Spenderabdeckung
und Anteil des größten beitragenden Spenders pro Cluster gespeichert.

Die Silhouetten sind insgesamt niedrig: 0,0793 für K-Means, 0,0671 für
Leiden und 0,0489 für Ward. Das unterstützt keine Behauptung perfekt
getrennter Zelltypen. K-Means liefert die stärkste kompakte Zusammenfassung
unter den getesteten Einstellungen; Leiden eine feinere Aufteilung mit
biologisch plausiblen zusätzlichen Profilen.
\source{Notebook 03, Abschnitte 3--4; \code{cluster_profiles};
\code{task3_donor_cluster_counts.csv}, \code{task3_cluster_counts.csv}.}

\clearpage
% Seite 12
\subsection*{Biologische Annotation: Evidenz statt Etikettensicherheit}
Die Annotation verwendet relative Markerprofile und die biologische
Einordnung der Ausgangsstudie\cite{horowitz}. „Kompatibel“ bezeichnet
Plausibilität, keine durch Referenzlabels bestätigte Identität. Insbesondere
CD117 allein definiert keine Stammzelle, CD3 mit NK-Rezeptoren allein keine
NKT-Zelle. Die IDs gelten nur innerhalb des jeweiligen Verfahrens.
\begin{table}[!ht]\centering\small
\caption{Vollständige Kurzannotation der 31 ausgewählten Cluster.
Die zugehörigen Markerwerte stehen in \code{task3_marker_profiles.csv}, die
ausführlichen Begründungen in \code{task3_annotations.csv}.}
\begin{tabular}{p{2.3cm}p{1cm}p{11.8cm}}\toprule
Verfahren & ID & Einordnung und entscheidende Markerhinweise\\\midrule
K-Means & 0 & B-kompatibel, breit: CD19/HLA-DR erhöht, mögliche Mischung\\
 & 1 & CD4-T-kompatibel: CD3/CD4/CD127\\
 & 2 & NK-kompatibel: CD56/CD16/CD94/NKp46; relativ niedrigeres CD3\\
 & 3 & CD8-T-kompatibel: CD3/CD8; wenig CD56/CD16\\
 & 4 & NK-kompatibel, KIR-reich: CD56/CD16 und 3DL1; CD3 intermediär\\
 & 5 & Myeloid-kompatibel, breit: CD33/HLA-DR/ILT2-CD85j\\\midrule
Ward & 0 & Gemischt, CD25-reich; zugleich CD19/CD33/HLA-DR\\
 & 1 & Myeloid-kompatibel: CD33/HLA-DR/ILT2-CD85j\\
 & 2 & CD4-T-kompatibel: CD3/CD4/CD127\\
 & 3 & B-kompatibel: CD19/HLA-DR\\
 & 4 & Unklar, NKp44-reich; kein eindeutiges NK-Gesamtprofil\\
 & 5 & CD8-T-kompatibel: CD3/CD8\\
 & 6 & Gemischt mit CD4-T-kompatiblem Anteil und CD19/CD33\\
 & 7 & Myeloid-kompatibel, CD117-reich\\
 & 8 & Gemischt: intermediäres CD3, CD57/NKp30\\
 & 9 & T-kompatibel mit NK-Rezeptoren: CD3/CD8, CD94/NKG2A\\
 & 10 & NK-kompatibel, NKG2A-reich: CD56/CD16/NKp46/CD94\\
 & 11 & NK-kompatibel, KIR-reich; intermediäres CD3\\
 & 12 & Unklar, TRAIL-reich\\
 & 13 & NK-kompatibel, CD57-reich; weniger NKG2A als Cluster 10\\\midrule
Leiden & 0 & Myeloid-kompatibel: CD33/HLA-DR/ILT2-CD85j\\
 & 1 & CD4-T-kompatibel: CD3/CD4/CD127\\
 & 2 & CD8-T-kompatibel: CD3/CD8\\
 & 3 & B-kompatibel: deutlicheres CD19/HLA-DR als K-Means 0\\
 & 4 & NK-kompatibel: CD56/CD16/CD94/NKp46\\
 & 5 & Unklar: schwache CD3/CD19/CD56/CD16-Linienmarker\\
 & 6 & T-kompatibel mit NK-Rezeptoren: CD3/CD8/CD57, CD94/NKG2A\\
 & 7 & Myeloid-kompatibel, CD117-reich\\
 & 8 & Gemischt, CD25-reich; intermediäres CD3/CD19/CD33\\
 & 9 & NK-kompatibel, KIR-reich; intermediäres CD3\\
 & 10 & Unklar, TRAIL-reich; CD3 messbar\\
\bottomrule\end{tabular}
\end{table}
\finding{Auf die Frage nach dem biologisch besten Verfahren lässt sich kein
eindeutiger Sieger belegen. K-Means gewinnt das festgelegte geometrische
Kriterium; Leiden ist eine nützliche feinere Beschreibung. Ohne Referenzlabels,
Coexpressionsprüfung aller Gruppen oder externe Annotation bleibt biologische
Richtigkeit offen. Dieser Unterschied ist Teil der Antwort auf Aufgabe 3.}

\clearpage
% Seite 13
\section{Aufgabe 4: gemeinsames Vergleichsprotokoll}\label{sec:protokoll}
\subsection*{Äußere Tests und innere Auswahl}
Wir übernehmen das spenderbezogene CMV-Schema des Papers: 100 Monte-Carlo-
Splits, jeweils sieben negative und sieben positive Trainingsspender.
Die übrigen vier negativen und zwei positiven Spender bilden den Test.
Die gespeicherte Splitdatei enthält daher \(100\times20=2.000\) Zeilen.
Sie ist die gemeinsame Quelle für alle drei Methoden.

\begin{figure}[!ht]\centering
\stepbox{13cm}{20 Spender: 11 CMV− und 9 CMV+\\Feste Label- und Markerzuordnung}
\par\(\downarrow\)\par
\stepbox{8cm}{14 äußere Trainingsspender: 7 CMV−, 7 CMV+\\Drei innere Folds zur Auswahl}
\hfill\stepbox{6cm}{6 äußere Testspender: 4 CMV−, 2 CMV+\\Bis zur fertigen Modellauswahl ausgeschlossen}
\par\(\downarrow\)\par
\stepbox{8cm}{Je innerem Fold: 9 oder 10 Trainingsspender\\4 oder 5 andere Spender zur Validierung}
\hfill\stepbox{6cm}{SVM / Citrus: finales 14-Spender-Modell\\CellCNN: ausgewähltes inneres Modell}
\caption{Erklärendes Schema des Informationsflusses, keine zusätzliche Analyse.
Alle Zellen eines Spenders folgen dessen Partition. Die rechte untere Box
beschreibt, welches fertige Modell auf die äußere Testgruppe angewandt wird.}
\label{fig:splits}
\end{figure}

Ein \code{default_rng(12345)} erzeugt für jeden Split einen uint32-Seed.
Mit diesem Seed werden pro Klasse sieben Trainingsspender ohne Zurücklegen
gewählt. Der Test ist exakt das Komplement. \code{StratifiedKFold(3,
shuffle=True, random_state=split_seed)} verteilt nur die 14 äußeren
Trainingsspender in innere Folds. Bei einem inneren Fold sind die anderen
beiden Folds Training. Äußere Testspender haben \code{inner_fold=-1}.

\textbf{Beispiel Split 0:} Der Seed ist 3.003.105.692. Jede Methode verwendet
die identische Spenderzuordnung. Der ausgewählte CellCNN-Kandidat stammt
aus innerem Fold 1 und ist auf neun Spendern trainiert; vier Filter werden
behalten. Die sechs äußeren Testspender sind weder an seinem Scaler-Fit
noch an Gewichtsoptimierung oder Kandidatenwahl beteiligt.

\textbf{Keine neue unabhängige Kohorte pro Wiederholung.} Ein Spender kann
in einem Split trainieren und in einem anderen testen. Innerhalb desselben
Splits wird er vollständig getrennt. Wiederholte Testauftritte sind abhängig;
600 Testvorhersagen pro Methode sind weiterhin Vorhersagen für nur 20 Personen.
\source{Notebook 04a, Abschnitt 9; \code{task4_donor_splits.csv};
CellCNN-Auswahl in Tabelle~\ref{tab:cnnbeispiel}.}

\clearpage
% Seite 14
\subsection*{Welche Daten darf welcher Schritt sehen?}
Die elementweise ArcSinh-Transformation lernt keine Parameter. Ihre
Berechnung vor der technischen Aufteilung ist deshalb bei festem Kofaktor
unproblematisch. Mittelwerte, Streuungen, Cluster und Regressionsparameter
sind dagegen gelernte Größen. Für die äußere Prüfung darf keiner dieser
Schritte die Testspender verwenden.

\begin{table}[!ht]\centering\small
\caption{Tatsächliche Datenmengen und Modellauswahl im Full-Modus.
„Alle“ meint alle Zellen des betreffenden Spenders im Alive-Gate.}
\begin{tabular}{p{4cm}p{3.6cm}p{3.6cm}p{3.6cm}}\toprule
Schritt & CellCNN & Citrus & SVM\\\midrule
Markerbasis & 37; ArcSinh, dann z & 37; ArcSinh & 37; ArcSinh, dann z\\
Zellen für Scaler & 20.000 je innerem Trainingsspender & Kein Marker-z; Standardisierung der Häufigkeitsfeatures im Fit & 10.000 je Trainingsspender\\
Trainingsinputs & 200 Gruppen je Spender mit je 3.000 Ziehungen & 1.000 Zellen je Spender & 10.000 Zellen je Spender\\
Zurücklegen & Ja, innerhalb der Inputs & Nein & Nein\\
Innere Auswahl & Beste Accuracy eines konkreten Kandidaten & Minimaler OOF-Klassifikationsfehler über Lambda & Mittlere Fold-ROC-AUC über C\\
Finale Trainingsspender & 9 oder 10; kein Refit & 14 & 14\\
Testauswertung & 5 Inputs à 20.000 Zellen & 1.000 gezogene Zellen & Alle\\
Spenderentscheidung & Mittlere Wahrscheinlichkeit ≥ 0,5 & Wahrscheinlichkeit ≥ 0,5 & Topscore ≥ OOF-Schwelle\\
\bottomrule\end{tabular}
\end{table}

\textbf{Spendergleichheit ist methodenspezifisch.} Gleiche Zellstichproben
geben Spendern im SVM-Zellverlust gleiches Gesamtgewicht. Gleiche Inputzahlen
bewirken dies im CellCNN-Trainingsverlust. Citrus hat eine Featurezeile je
Spender. Bei Auswertungen liefert jeder Testspender genau einen Score.
Zwei positive und vier negative Testspender bedeuten trotzdem nicht gleiche
Klassenhäufigkeit; Balanced Accuracy behandelt die beiden Klassen separat.

\textbf{Grenze der inneren Trennung bei Citrus.} Die originalen Funktionen
erzeugen den gemeinsamen Lambda-Pfad aus den Features und Labels aller
14 äußeren Trainingsspender. Die foldweisen Hierarchien und Fits verwenden
ihre jeweiligen inneren Trainingsspender, aber das Kandidatenraster kennt
indirekt die inneren Validierungslabels. Eine vollständig isolierte innere
Suche wäre strenger. Diese Präzisierung betrifft den originalen
Auswahlmechanismus; ein Zugriff auf äußere Testspender wurde dabei nicht
gefunden (Details S.~\pageref{sec:citruslambda}).

\textbf{Fortsetzen und Nachweise.} Methodenläufe speichern Vorhersagen,
Auswahl und Modellparameter nach jedem äußeren Split. Konfigurationen
enthalten Parameter, Versionen, Code und Eingabeprüfsummen. Abweichende oder
unvollständige Zustände dürfen nicht als passende Checkpoints fortgesetzt
werden. Der Bericht lädt nur die vorhandenen vollständigen Full-Ergebnisse.
\source{\code{src/task4_artifacts.py}, \code{src/task4_artifacts.R};
Methoden-Notebooks und Konfigurationsdateien mit Suffix \code{gated_alive_full}.}

\clearpage
% Seite 15
\section{CellCNN: vom Zellprofil zur Spenderwahrscheinlichkeit}\label{sec:cellcnn}
\subsection*{Architektur und Bedeutung der Tensoren}
CellCNN betrachtet eine Probe als ungeordnete Menge. Derselbe lineare
Filter wird auf jede Zelle angewandt, unabhängig von ihrer Position.
Es gibt keine Nachbarschaft zwischen aufeinanderfolgenden FCS-Zeilen,
die wie ein Bildpixelmuster gefaltet werden müsste. \code{nn.Linear}
implementiert hier den relevanten gemeinsamen Filteroperator.

Ein Batch enthält \(B\) Multi-Cell-Inputs mit je \(n\) Zellen und \(m=37\)
Markern. Der Tensor hat Form \(B\times n\times37\). Bei \(K\) Filtern
entstehen \(B\times n\times K\) Antworten. Für standardisierte Zelle
\(z_i\) berechnet Filter \(k\)
\[
 a_{ik}=b_k+\sum_{j=1}^{37}W_{kj}z_{ij},\qquad
 r_{ik}=\max(0,a_{ik}).
\]
\(W\) sind Gewichte, \(b\) Biases. ReLU setzt negative lineare Antworten
auf null. Ein großer positiver Wert bedeutet gute Übereinstimmung mit
\emph{diesem gelernten linearen Muster}, nicht automatisch hohe Expression
aller Marker oder eine positive CMV-Klasse.

\begin{figure}[!ht]\centering
\stepbox{14cm}{Input: \(B\times3.000\times37\) standardisierte Markerwerte}
\par\(\downarrow\)\par
\stepbox{14cm}{Gemeinsame lineare Filter und ReLU: \(B\times3.000\times K\)}
\par\(\downarrow\)\par
\stepbox{14cm}{Je Filter Mittel der 30 höchsten Antworten: \(B\times K\)}
\par\(\downarrow\)\par
\stepbox{14cm}{Lineare Ausgabe: \(B\times2\) Logits; Softmax liefert zwei Wahrscheinlichkeiten}
\caption{Erklärendes Architekturschema für Trainingsinputs. Im Full-Modus
ist \(B\le128\), \(K\in\{3,4,5\}\). Bei Testinputs mit 20.000 Zellen
werden 200 Antworten je Filter gepoolt.}\label{fig:cnn}
\end{figure}

Die Zahl \(K\) ist klein: Zusammen einschließlich beider Biasvektoren hat
das Netz \(37K+K+2K+2=40K+2\) Parameter, also 122, 162 oder 202.
Es gibt keine zusätzliche versteckte Schicht und kein Dropout.

Permutation der Zellen ändert den Vorwärtsscore mathematisch nicht:
Filter arbeiten zeilenweise, Top-k-Pooling ignoriert die Reihenfolge.
Distanzkarten aus Aufgabe 2 werden dem Netz nicht zugeführt. Das Netz
sieht die 37 Marker direkt und lernt die Zellrepräsentation über den
Spenderverlust. Diese Eigenschaft begründet den Vergleich mit einer
separaten unüberwachten Clusterbildung.
\source{Notebook 04c, Abschnitt 3, Klasse \code{CellCNN}; Parameterexport
in \code{task4_cellcnn_filters_gated_alive_full.csv}.}

\clearpage
% Seite 16
\subsection*{Trainings-Scaler und Erzeugung der Multi-Cell-Inputs}
Für jeden Kandidaten werden ausschließlich dessen neun oder zehn innere
Trainingsspender verwendet, um den Scaler zu fitten. Aus jedem werden
20.000 Zellen ohne Zurücklegen gezogen. Die verkettete Matrix hat somit
180.000 oder 200.000 Zeilen und 37 Spalten. \(\mu_j\) und \(\sigma_j\)
werden aus dieser Matrix berechnet und unverändert auf Training,
Validierung und später Test angewandt. Ein neuer Scaler für jeden
Testspender würde das gelernte Koordinatensystem verändern und ist nicht
Teil der Methode.

Dann erzeugt jeder Trainingsspender 200 Inputs mit je 3.000 Zellziehungen
\emph{mit Zurücklegen}. Ein Kandidat erhält damit 1.800 oder 2.000
Trainingsinputs. Innerhalb eines Inputs dürfen dieselben Originalzellen
mehrfach vorkommen; zwischen Inputs ebenso. Die 200 Inputs sind zusätzliche
Trainingsbeispiele, aber keine 200 unabhängigen Menschen.

\textbf{Warum große zufällige Zellmengen?} Ein Input soll mit ausreichender
Wahrscheinlichkeit Zellen einer seltenen Population enthalten. Für einen
hypothetischen Anteil \(p\) beträgt bei unabhängigen Ziehungen mit Zurücklegen
die Wahrscheinlichkeit mindestens einer solchen Zelle
\[
 P(\text{mindestens eine})=1-(1-p)^{3000}.
\]
Für \(p=0{,}001\) ergibt das ungefähr 95,0\%, für \(p=0{,}0001\) etwa
25,9\%. Dies ist eine illustrative Samplingrechnung, keine Schätzung des
wahren CMV-Subsets in unseren Daten. Sie zeigt auch, weshalb 3.000 Zellen
nicht jede beliebig seltene Population zuverlässig erfassen.

Die vier oder fünf Validierungsspender liefern ebenfalls je 200 feste
3.000-Zell-Inputs, also 800 oder 1.000 Beispiele für den Validierungsverlust.
Sie werden mit dem Trainings-Scaler transformiert. Ihre Zellen werden
nicht zum Gradientenupdate verwendet.

\textbf{Feste Inputs, wechselnde Batchreihenfolge.} Die Funktion
\code{materialize_multicell_inputs} erzeugt sämtliche Inputs einmal je
Kandidat und hält sie als Tensoren auf dem Rechengerät. Die Zellen werden
nicht in jeder Epoche neu gezogen. Der Trainings-DataLoader mischt die
Reihenfolge dieser festen Inputs; Validierungsinputs bleiben in fester
Reihenfolge. Batchgröße ist 128; der letzte Batch darf kleiner sein.

\textbf{Seed-Kette.} Für äußeren Seed \(s\), inneren Fold \(f\) und
Filterzahl \(K\) lautet der Kandidatenseed \(s+10.000f+100K\).
Er steuert Gewichtsinitialisierung und Scaler-Sampling. Trainingsinputs
beginnen mit Offset 1.000, Validierungsinputs mit 2.000; jeder Input addiert
seine laufende Indexposition. Die Konstruktion sortiert Spender-IDs.
Da der Kandidatenseed auch \(K\) enthält, unterscheiden sich zwischen
Filterzahlen nicht nur Gewichte, sondern auch zufällige Stichproben.

Python-Modelldaten werden zuerst nach Float32 konvertiert, danach wird
ArcSinh berechnet. Diese Reihenfolge ist in der späteren Modellrekonstruktion
beizubehalten; Float64-ArcSinh mit anschließendem Cast ist nicht bitgleich.
\source{\code{fit_balanced_scaler}, \code{scale_donors},
\code{materialize_multicell_inputs}, \code{train_candidate}.}

\clearpage
% Seite 17
\subsection*{Pooling, Ausgabe und Optimierung}
Für Filter \(k\) seien die Antworten eines Inputs absteigend sortiert.
Die Implementierung definiert
\[
 q(n)=\max\{1,\lfloor0{,}01n\rfloor\},\qquad
 h_k=\frac1{q(n)}\sum_{\ell=1}^{q(n)}r_{(\ell)k}.
\]
Es wird \emph{abgerundet}: bei 3.000 Zellen sind es 30, bei 20.000 sind
es 200. Verschiedene Filter dürfen verschiedene Zellen auswählen. Auch
Nullantworten können in die Top-k-Menge gelangen, falls zu wenige positive
Antworten vorhanden sind. Pooling misst eine gewichtete Reaktionsstärke;
es ist weder eine direkt gezählte Zellhäufigkeit noch ein harter Zelltypanteil.

Die Ausgabeschicht bildet aus \(h\) zwei Logits:
\[
 \ell_c=d_c+\sum_k V_{ck}h_k,\qquad
 P(c\mid X)=\frac{e^{\ell_c}}{e^{\ell_0}+e^{\ell_1}}.
\]
Für CMV+ ist die Differenz \(\ell_1-\ell_0\) entscheidend. Daher hat
Filter \(k\) die Richtung \(\Delta V_k=V_{1k}-V_{0k}\). Erst dieser
Kontrast bestimmt, ob mehr Filterantwort das Netz in Richtung CMV+ oder
CMV− bewegt.

Die Verlustfunktion pro Batch ist
\[
 L=-\frac1B\sum_{b=1}^B\log P(y_b\mid X_b)
   +10^{-4}\left(\sum_{kj}W_{kj}^2+\sum_{ck}V_{ck}^2\right).
\]
Beide Gewichtsmatrizen werden L2-bestraft, \emph{beide Biasvektoren nicht}.
Die Strafe wird explizit addiert, nicht als AdamW-Gewichtszerfall umgesetzt.
\code{cross_entropy} erhält Logits und berechnet die nötige Normalisierung
intern; vor dem Verlust wird kein zusätzliches Softmax angewandt.

Adam aktualisiert nach \code{zero_grad}, Vorwärtsrechnung,
\code{backward} und \code{step} alle Netzparameter. Lernrate ist 0,01;
die unveränderten PyTorch-Adam-Defaults sind \(\beta_1=0{,}9\),
\(\beta_2=0{,}999\), \(\epsilon=10^{-8}\), ohne Weight Decay.
\code{nn.Linear} verwendet seine Standardinitialisierung nach dem
Kandidatenseed, keine aus dem alten Keras-Netz importierten Gewichte.

\textbf{Early Stopping.} Nach jeder Epoche wird ohne Gradienten der
regularisierte Verlust aller festen Validierungsinputs berechnet.
Batchverluste werden nach ihrer Beispielzahl gewichtet. Ein Zustand wird
nur gespeichert, wenn der Verlust um mehr als \(10^{-6}\) sinkt.
Nach fünf Epochen ohne solche Verbesserung stoppt der Kandidat, spätestens
nach 100 Epochen. Anschließend wird der beste gespeicherte Zustand geladen,
nicht einfach der letzte.

\finding{Der historische Trainingsverlust ist ein ungewichtetes Mittel der
Batchmittel, der Validierungsverlust hingegen beispielgewichtet. Der
Trainingsmittelwert steuert den Stopp nicht. Epochenkurven werden nicht
nachträglich erfunden: Die vollständige Historie wird vom Training erzeugt,
aber nicht in den dauerhaften CSV-Artefakten exportiert.}
\source{\code{regularized_loss}, \code{evaluate_loader}, \code{train_candidate};
PyTorch-Quellcode der installierten Version für Optimiererdefaults.}

\clearpage
% Seite 18
\subsection*{Neun Kandidaten, eine explizite Auswahlregel}
Je äußerem Split werden drei Filterzahlen in drei inneren Folds trainiert.
Das ergibt neun Kandidaten und insgesamt 900 Fits im vorhandenen
Full-Benchmark. Es wird nicht zuerst die mittlere Leistung einer Filterzahl
über die Folds ermittelt. Stattdessen konkurrieren neun konkrete Netze,
die auf unterschiedlichen Trainings-/Validierungsmengen entstanden sind.

Für die Kandidatenbewertung liefert jeder innere Validierungsspender eine
Wahrscheinlichkeit: fünf zufällige Inputs mit bis zu 20.000 Zellen,
Vorwärtsrechnung, dann Mittelwert. Daraus werden Accuracy bei Schwelle 0,5
und ROC-AUC auf \emph{Spenderebene} berechnet. Das ist eine andere Bewertung
als der Verlust der 3.000-Zell-Inputs, der Early Stopping steuert.

Die deterministische Sortierung verwendet der Reihe nach:
\begin{enumerate}
\item höchste Validierungs-Accuracy;
\item bei Gleichstand höchste Validierungs-ROC-AUC;
\item danach kleinster bester Validierungsverlust;
\item danach weniger Filter;
\item zuletzt kleinere innere Foldnummer.
\end{enumerate}
\tab{cellcnn_split0}{Tatsächliche neun Kandidaten für äußeren Split 0.
Accuracy und AUC beziehen sich auf vier bzw. fünf Validierungsspender.
„Epochen“ ist die Zahl durchlaufener Epochen, nicht die beste Epoche.}{tab:cnnbeispiel}{}

Der Kandidat mit Fold 1 und vier Filtern gewinnt mit Accuracy 0,8. Der
Kandidat mit Fold 2 und fünf Filtern erreicht zwar AUC 1,0, aber nur
Accuracy 0,75 und wird nicht gewählt. Diese Tabelle macht prüfbar,
welches Auswahlziel tatsächlich Priorität hat. Primäre \emph{äußere}
Berichtsmetrik ROC-AUC und primäres \emph{inneres} CellCNN-Auswahlkriterium
Accuracy sind unterschiedliche Entscheidungen.

\tab{cellcnn_auswahl}{Verteilung der 100 ausgewählten Netze nach Filterzahl
und innerem Fold. Quelle: alle 900 Kandidatenzeilen, Auswahlflag und
Sortierschlüssel unabhängig abgeglichen.}{tab:cnnauswahl}{}

\finding{Die ausgewählten Netze haben zusammen \CellcnnFilters{} Filter.
Sie wurden im Median \CnnEpochMedian{} Epochen trainiert; \CnnMaxEpochs{}
erreichen die Grenze von 100. Das sagt nichts über die Lage der besten
Epoche aus. Die häufigere Auswahl eines Folds oder einer Filterzahl ist
keine unabhängige biologische Bestätigung.}
\source{\code{train_outer_split}; \code{task4_cellcnn_selection_gated_alive_full.csv}.}

\clearpage
% Seite 19
\subsection*{Endgültige Vorhersage und gespeichertes Modell}
Nach der Auswahl wird das Netz \textbf{nicht auf 14 Spendern neu trainiert}.
Es behält seinen besten Zustand und den Scaler der neun oder zehn inneren
Trainingsspender. Dies orientiert sich am im Paper beschriebenen NK-
Auswahlschema. Die übrigen inneren Validierungsspender haben Auswahl und
Stopp beeinflusst, aber nicht seine Gradienten oder seinen Scaler-Fit.

Für einen äußeren Testspender \(d\) werden fünf Inputs \(X_d^{(1)},\dots,
X_d^{(5)}\) gezogen. Jeder enthält \(\min(20.000,N_d)\) Zellen ohne
Zurücklegen. Zwischen den fünf Inputs können Zellen mehrfach vorkommen.
Es gibt keine gemeinsame feste Partition in fünf disjunkte Mengen.
Die Vorhersage lautet
\[
 \widehat p_d=\frac15\sum_{b=1}^5P(1\mid X_d^{(b)}),\qquad
 \widehat y_d=\mathbf1\{\widehat p_d\ge0{,}5\}.
\]
Dies ist der Mittelwert der Wahrscheinlichkeiten, nicht das Softmax eines
mittleren Logits. Beispiel: Die fünf illustrativen Werte 0,4; 0,6; 0,7;
0,8; 0,9 ergeben \(\widehat p_d=0{,}68\) und Klasse 1.

Der Zufallsgenerator wird mit \(s+900.000+i\) initialisiert, wobei \(s\)
der äußere Splitseed und \(i\) die Position des Spenders in der sortierten
Testliste ist. Fünf aufeinanderfolgende Ziehungen mit demselben Generator
machen den Samplingablauf rekonstruierbar. Bei der inneren Kandidatenbewertung
lautet der entsprechende Offset zum Kandidatenseed 50.000.

\begin{table}[!ht]\centering\small
\caption{Welche Parameter den ausgewählten CellCNN vollständig beschreiben.}
\begin{tabular}{p{5cm}p{10cm}}\toprule
Exportinhalt & Bedeutung für die Wiederverwendung\\\midrule
Filtergewicht pro Marker & Matrix \(W\), gleiche Markerreihenfolge wie beim Fit\\
Filterbias & Vektor \(b\), je Filter in den Markerzeilen wiederholt\\
Beide Ausgabegewichte & \(V_{0k}\), \(V_{1k}\); ihr Kontrast bestimmt die Richtung\\
Beide Ausgabebiases & \(d_0,d_1\); notwendig für vollständige Wahrscheinlichkeiten\\
Scaler-Mittel und -Skala & \(\mu_j,\sigma_j\) des ausgewählten Kandidaten\\
Fold, Seed, Gate, Gerät & Zuordnung zu Trainingsspendern und Rechenbedingungen\\
Poolinganteil und Kofaktor & Rechenregeln für Antwortaggregation und Transformation\\
\bottomrule\end{tabular}
\end{table}

CSV-Dateien werden für die Rekonstruktion mit
\code{float_precision="round_trip"} gelesen. Aufgabe 5 stellt einen
\code{StandardScaler} wieder her, statt dessen Transformationsarithmetik
durch eine vermeintlich äquivalente Formel mit anderer Gleitkommagenauigkeit
zu ersetzen. Auch das verwendete Rechengerät wird dokumentiert: Der
vorhandene vollständige CellCNN-Lauf verwendete CUDA.

Die gespeicherten Aufgabe-5-Checks enthalten 600 rekonstruierte
Spenderwahrscheinlichkeiten für CellCNN und zugehörige Logitidentitäten.
Ihre aktuelle Dateiidentität ist geprüft; für diesen Bericht werden
nicht nochmals 100 Netze trainiert oder alle Inferenzläufe wiederholt.
\source{\code{predict_donor}; Export in \code{train_outer_split};
\code{SavedCellCNN} und \code{restore_scaler} in Aufgabe 5.}

\clearpage
% Seite 20
\subsection*{Paper-Nähe, Referenznotebooks und bewusst andere Entscheidungen}
Die beiden offiziellen Notebooks wurden beim Bericht lesend geprüft.
Das NK-Beispiel verwendet \code{ncell=200,nsubset=1000}; das Alive-Beispiel
\code{ncell=3000,nsubset=1000,maxpool_percentages=[0.5,1]}.
Beide demonstrieren einen einzelnen zufälligen äußeren Split mit einer
weiteren festen inneren Teilung und rufen die historische Bibliothek auf.
Sie sind daher nicht identisch mit dem 100-Split-Protokoll des Papers.

\begin{table}[!ht]\centering\small
\caption{Abgleich des Projekts mit dem CMV-Protokoll des Papers bzw. der
Referenzimplementierung. Unterschiede sind keine stillschweigenden Gleichsetzungen.}
\begin{tabular}{p{3.5cm}p{5.5cm}p{6cm}}\toprule
Aspekt & Bezug & Umsetzung und Bedeutung\\\midrule
Äußere Auswertung & Paper: 100 Splits, 7+7 Training & Übernommen; identische Splits für alle Methoden\\
Innere Auswahl & Paper: 3 Folds, Zufallssuche & 3 Folds, festes Raster von 3/4/5 Filtern; kleinere Suche\\
Inputs & Paper: 3.000 Zellen, 200 Inputs je Probe & Übernommen; feste materialisierte Inputs je Kandidat\\
Pooling & Paper: Mittel der höchsten 1\% & Übernommen mit genauer Abrundungsregel\\
Optimierung & Paper: Lernrate 0,01, kein Dropout & Übernommen; aktuelle PyTorch-Initialisierung/Adam statt altem Framework\\
Skalierung & Gelernte Trainingsnormalisierung & Gleich viele Scaler-Zellen je Trainingsspender; andere Stichproben als Legacy-Code\\
Kandidatenbewertung & Höchste Validierungsgenauigkeit & Spenderweise Mittelung; zusätzliche festgelegte Tie-Breaker\\
Finales Netz & Paper: bestes Netz aus inneren Läufen & Kein Refit auf allen äußeren Trainingsspendern\\
Testinputs & Paper: ein Input mit 20.000 Zellen & Fünf Inputs; geringere Samplingstreuung beabsichtigt, nicht separat als Verbesserung nachgewiesen\\
Interpretation & Paper: Halbmaximum-Subsets und Gruppen ihrer Zentroiden & Trainingsgebundene Halbmaxima, OOF-Zellhäufigkeiten; keine Zentroidgruppenreplikation\\
Markerzahl & Papertext 36, Beispiele 37 & Alle 37 gelieferten Marker\\
\bottomrule\end{tabular}
\end{table}

\textbf{Was ist damit reproduziert?} Wir bilden den biologischen
Problemtyp, die zentrale Architektur und das spenderbezogene CMV-
Evaluationsschema nach. Wir behaupten keine numerisch exakte Replikation
aller Paperresultate, keine Reproduktion der alten Zufallssuche und keine
identische Interpretation der wiederkehrenden Zellpopulationen.

\textbf{Was fehlt für stärkere Aussagen?} Eine isolierte Untersuchung des
Nutzens der fünf Testinputs, ein Vergleich mit einem Refit auf 14 Spendern
oder eine systematische Ablation der Filterantwort sind nicht vorhanden.
Diese Änderungen würden neue, ausdrücklich getrennte Experimente erfordern.
Die optionale Aufgabe 6 wurde bislang nicht ausgeführt.
\source{Paper, S. 8, „NK-cell benchmark data set“; offizielle Beispiele
\cite{cellcnncode}; \code{task4_cellcnn_predictions_gated_alive_full.config.json}.}

\clearpage
% Seite 21
\section{Citrus: Hierarchie, Häufigkeiten und Regression}\label{sec:citrus}
\subsection*{Originalimplementation und Zellraum}
Citrus\cite{citrus} wird in einer separaten R-Umgebung verwendet, nicht als
vereinfachte Python-Ersatzmethode. Fixiert sind Citrus 0.8, Commit
\code{d02baae544abdc403704aaceb75d1e7931a0331c}, und Rclusterpp,
Commit \code{a07380683ce7a6849af8ec27db6439ea3a707890}.
Die Funktionen und Defaults der installierten Pakete wurden für diesen
Bericht lesend kontrolliert\cite{citruscode}.

Pro Spender zieht Citrus 1.000 Zellen ohne Zurücklegen, nach der
festen Reihenfolge der sortierten Dateiliste. Der äußere Trainingsbaum
umfasst 14.000 Zellen. Marker werden mit ArcSinh und Kofaktor 5
transformiert; \textbf{kein zusätzliches Marker-z-Scoring} wird vor dem
Clustering angewandt. Dies unterscheidet Citrus von Aufgabe 3, deren Ward-
Clustering auf der standardisierten explorativen Stichprobe arbeitet.

\code{citrus.cluster.hierarchical} ruft \code{Rclusterpp.hclust} auf.
Dessen hier geprüfte Defaults sind Ward-Linkage und euklidische Distanz.
Die Cluster sind Knoten einer Hierarchie: Eltern enthalten die Zellen ihrer
Kinder. Sie sind keine einzige disjunkte Partition mit festgelegter
Clusterzahl. Eine Zelle kann mehreren berücksichtigten Knoten angehören.

\begin{figure}[!ht]\centering
\stepbox{14cm}{14 Trainingsspender × 1.000 Zellen; 37 ArcSinh-Marker}
\par\(\downarrow\)\par
\stepbox{14cm}{Ward-Hierarchie; Knoten mit mindestens 5\% der Trainingszellen}
\par\(\downarrow\)\par
\stepbox{14cm}{Pro Spender eine Zeile: relative Häufigkeit jedes zulässigen Knotens}
\par\(\downarrow\)\par
\stepbox{14cm}{L1-logistische Regression; ausgewählte Clusterkoeffizienten}
\caption{Erklärendes Citrus-Schema. Die entsprechende Hierarchie wird auch
innerhalb jedes inneren Folds nur aus dessen Trainingsspendern aufgebaut.
Überlappende Cluster sind ausdrücklich Teil der Repräsentation.}\label{fig:citrus}
\end{figure}

\code{minimumClusterSizePercent=0.05} bedeutet einen Anteil von 5\%,
nicht 0,05\%. Im äußeren Baum entspricht das mindestens 700 Zellen;
bei 9.000 bzw. 10.000 inneren Trainingszellen sind es 450 bzw. 500.
Ein seltenes Subset unterhalb dieser Grenze kann keinen so kleinen
eigenständigen Cluster als Merkmal liefern, möglicherweise aber größere
Clusterhäufigkeiten verändern. Die Grenze ist daher eine plausible
Sensitivitätsbeschränkung, kein Beweis einer zwangsläufigen Fehlklassifikation.
\source{Notebook 04d, \code{read_citrus_files}, \code{build_fixed_fold_clustering};
\code{citrus.cluster.hierarchical}, Defaults von \code{Rclusterpp.hclust}.}

\clearpage
% Seite 22
\subsection*{Von der Hierarchie zu einer Spender-Featurematrix}
Für Cluster \(k\) mit Zellmenge \(C_k\) und die gezogenen Zellen
\(I_d\) eines Spenders ist das Häufigkeitsmerkmal
\[
 F_{dk}=\frac{|I_d\cap C_k|}{|I_d|}.
\]
Im finalen Training ist \(|I_d|=1.000\). Ein Beispielcluster mit
120 Zellen eines Spenders liefert somit \(F_{dk}=0{,}12\), nicht 12.
Das R-Paket verwendet Anteile zwischen null und eins. Fehlende
Mitgliedschaft ergibt null. Weil die Cluster überlappen, müssen ihre
Spaltenwerte pro Spender nicht zu eins summieren.

Die Featurematrix hat 14 Zeilen und eine vom Trainingsbaum abhängige
Spaltenzahl. Die Labels gehören in derselben Reihenfolge zu diesen
Spenderzeilen. Der Faktor hat die Klassenreihenfolge \code{CMV-},
\code{CMV+}; die ausgegebene positive Wahrscheinlichkeit ist CMV+.

\textbf{Logistischer Prädiktor.} Auf Häufigkeitsmerkmalen lautet das Modell
\[
 \eta_d=\beta_0+\sum_k\beta_kF_{dk},\qquad
 p_d=\operatorname{sigmoid}(\eta_d)=\frac1{1+e^{-\eta_d}}.
\]
Eine positive \(\beta_k\) erhöht bei sonst gleichen Features den Logit,
eine negative senkt ihn. Die Hierarchie erzeugt jedoch korrelierte,
überlappende Merkmale. Ein einzelner Koeffizient ist daher nicht die
unabhängige kausale Wirkung eines Zelltyps.

Mit \code{glmnet}, binomialer Familie und \(\alpha=1\) wird sinngemäß
\[
 \min_{\beta_0,\beta}\left\{-\frac1D\sum_{d=1}^D
 [y_d\log p_d+(1-y_d)\log(1-p_d)]
 +\lambda\sum_k|\widetilde\beta_k|\right\}
\]
optimiert. Die Tilde erinnert daran, dass die Regularisierung auf intern
standardisierten Features wirkt. \code{standardize=TRUE} zentriert und
skaliert die \emph{Clusterhäufigkeiten} im jeweiligen Trainingsfit;
der Intercept wird nicht penalisiert. Die exportierten Koeffizienten
werden von glmnet wieder auf die ursprüngliche Feature-Skala zurückgegeben.
Sie können deshalb später direkt mit den unstandardisierten Häufigkeiten
multipliziert werden\cite{glmnet}.

\textbf{Was bewirkt L1?} Bei größerem \(\lambda\) wird eine sparsamere
Lösung bevorzugt; viele Koeffizienten werden exakt null. Das liefert die
Auswahl der mit dem Label verbundenen Cluster. Die Größe des Rasters
soll nicht mit der Zahl unabhängiger Beobachtungen verwechselt werden:
Auch Hunderte Features werden aus lediglich 14 Spenderzeilen geschätzt.

Zusätzlich exportiert das Notebook für jeden wirksamen Cluster den
\emph{arithmetischen Mittelwert} jedes transformierten Markers über dessen
Trainingszellen. Dieser Zentroid ist eine Beschreibung des ausgewählten
Knotens. Er ist weder ein trainierter Markerkoeffizient noch identisch mit
den spendergewichteten Medianprofilen aus Aufgabe 5.
\source{\code{citrus.calculateFileClusterAbundance},
\code{citrus.buildModel.classification}, \code{extract_selected_cluster_profiles}.}

\clearpage
% Seite 23
\subsection*{Innere Auswahl und der gemeinsame Lambda-Pfad}\label{sec:citruslambda}
Die öffentliche Citrus-Funktion zur Foldbildung würde eigene zufällige
Folds erzeugen. Deshalb baut das Notebook das benötigte
\code{citrus.foldClustering}-Objekt mit den bereits gespeicherten
Spenderfolds auf. Dazu werden originale Funktionen zum Fold-Clustering
und Mapping verwendet. Clustering, Häufigkeitsberechnung und Regression
selbst werden nicht durch neue Verfahren ersetzt.

Für jeden der drei inneren Folds entstehen ein Baum aus den neun oder zehn
Trainingsspendern, Features dieser Trainingsspender und gemappte Features
der vier oder fünf Validierungsspender. Die Validierungszellen werden in
die feste Trainingshierarchie eingeordnet. Ihre Marker werden nicht für
den inneren Baumfit verwendet. Spalten sind foldinterne Cluster; dieselbe
numerische Cluster-ID in verschiedenen Bäumen hat keine gemeinsame Identität.

\textbf{Regulierung auswählen.} Die Originalfunktion
\code{citrus.endpointRegress} erzeugt zunächst einen absteigenden
Lambda-Pfad aus \code{allFeatures} und den Labels aller 14 äußeren
Trainingsspender. \code{nlambda=100} ist die angeforderte Pfadlänge,
nicht zwingend die Anzahl tatsächlich zurückgelieferter Lambda-Werte.
An denselben Lambda-Werten werden die inneren Fits bewertet. Die
OOF-Klassifikationsfehler der 14 Spender werden zusammengeführt.
Gewählt wird \code{cv.min}: minimale Fehlerrate, bei Gleichstand der
erste Eintrag im absteigenden Pfad, also stärkere Regularisierung.
Die zusätzlich berechnete 1-SE-Alternative wird nicht verwendet.

\finding{Das Lambda-Raster ist datenabhängig und sieht alle 14 äußeren
Trainingslabels. Damit ist es nicht unabhängig von den inneren
Validierungsspendern. Dies ist im Originalpaket so implementiert und
wird hier übernommen. Der Satz „jede datenabhängige Entscheidung wird nur
auf inneren Trainingsspendern gefittet“ wäre für Citrus zu weitgehend.
Die sechs äußeren Testspender bleiben weiterhin ausgeschlossen; die
berichtete äußere Leistung bewertet genau dieses Auswahlverfahren.}

Das endgültige Regressionsmodell wird auf den Häufigkeitsfeatures des
14-Spender-Baums trainiert. Technisch berechnet das Paket den finalen
Regularisierungspfad bereits während \code{endpointRegress}; die innere
CV bestimmt anschließend, welcher Punkt dieses Pfads für Vorhersagen
verwendet wird. „Nach Auswahl refitten“ beschreibt daher die statistische
Rolle, nicht exakt die Reihenfolge jedes internen Funktionsaufrufs.

Der auf R abgebildete Seed ist
\(s_R=s\bmod(2^{31}-1)\). Die Modulo-Abbildung verändert weder die
Spenderzuordnung noch die gemeinsamen Folds. Sie passt den uint32-Seed
an Rs ganzzahligen Wertebereich an. Training verwendet \(s_R\),
das äußere Testsampling \(s_R+1\).

\textbf{Nachweisgrenze.} Die finalen CSVs speichern das gewählte Lambda,
Clusterzahl und Vorhersagen, aber nicht sämtliche inneren
Lambda-Fehlerraten. Deren vollständige historische Auswahl lässt sich
allein aus diesen CSVs nicht neu nachrechnen. Ältere gezielte Prüfungen
existieren; der Bericht trainiert dafür keine neuen Citrus-Modelle.
\source{Jetzt lesend geprüft: \code{citrus.endpointRegress},
\code{citrus.generateRegularizationThresholds.classification},
\code{citrus.getCVMinima}; Notebook 04d, Abschnitt 4.}

\clearpage
% Seite 24
\subsection*{Neue Spender zuordnen und Nullmodelle verstehen}
Jeder äußere Testspender liefert 1.000 neue gezogene Zellen. Für jede
Testzelle wird eine nächste Zelle des festen Trainingsdatensatzes in den
37 ArcSinh-Markern gesucht. Die Testzelle übernimmt deren Mitgliedschaften
in der Hierarchie. Es wird weder ein Baum mit Testzellen ergänzt noch
lediglich der nächste Clusterzentroid gewählt. Die resultierenden
Häufigkeiten werden in exakt derselben Spaltenreihenfolge wie die
Trainingsfeatures an die finale Regression übergeben.

\textbf{Illustratives Mapping.} Eine nächste Trainingszelle gehört zum
Elternknoten A und dessen Kind B. Dann erhöht die Testzelle beide
Zählungen. Bei \(\beta_A=2\) und \(\beta_B=-3\) beträgt ihr addierter
Zellbeitrag \(-1\). Das Beispiel erklärt, warum „liegt in einem positiven
Cluster“ bei überlappenden Gruppen nicht für eine positive Nettorichtung
ausreicht.

Vor dem Export werden die Testwahrscheinlichkeiten auf 15 Nachkommastellen
gerundet. Wirksame Cluster werden konsistent durch
\(|\beta_k|>10^{-10}\) definiert. Die Toleranz betrifft Zählung und
Interpretation; sie ist von der Rundung der Wahrscheinlichkeiten zu
unterscheiden. Werte in der Größenordnung \(10^{-15}\) werden nicht als
biologisch relevante Clusterkoeffizienten ausgegeben.
\fig{modellauswahl}{1}{Vorhandene Modellauswahl im Hauptlauf. Links
Epochenzahlen aller 900 CellCNN-Kandidaten und der 100 ausgewählten Netze;
keine Lernkurven. Rechts Zahl wirksamer Citrus-Cluster pro äußerem Split
nach \(|\beta|>10^{-10}\). Die 15 Nullmodelle bleiben Teil des Vergleichs.}{fig:auswahl}

In 15 Splits sind keine wirksamen Cluster vorhanden. Bei sieben positiven
und sieben negativen Trainingsspendern liefert das reine Interceptmodell
\(p=0{,}5\). Mit der Regel \(p\ge0{,}5\) werden alle Testspender positiv
klassifiziert. Dann sind Sensitivität 1, Spezifität 0, Balanced Accuracy
0,5 und ROC-AUC 0,5. Dies ist eine gültige, aber uninformative Modelllösung.
Die Nullmodelle werden weder im Leistungsbericht noch aus den Nennern
von Aufgabe 5 entfernt.

\finding{Schwächere Citrus-Ergebnisse unter diesen Einstellungen beweisen
nicht, dass Citrus grundsätzlich ungeeignet ist. 1.000 Zellen je Spender,
5\%-Mindestclustergröße, hoch korrelierte Cluster und kleine innere Folds
beschränken die konkrete Analyse. Deren getrennte Effekte wurden nicht
systematisch abgetragen.}
\source{\code{run_citrus_split}, \code{citrus.mapToClusterSpace};
\code{task4_citrus_selection_gated_alive_full.csv}; \code{src/task4_artifacts.R}.}

\clearpage
% Seite 25
\section{Lineare Single-Cell-SVM im Detail}\label{sec:svm}
\subsection*{Die Modellannahme und der Zellverlust}
Jede Trainingszelle erhält das binäre Label ihres Spenders. Das ist eine
\emph{schwache Beschriftung}: Nicht jede Zelle eines CMV+-Spenders muss
selbst eine krankheitsassoziierte Zelle sein. Die SVM erhält aber keine
feineren Trainingslabels. Diese Annahme stammt aus der Zellbaseline des
Papers; unsere spenderweise Aggregation wird zusätzlich festgelegt.

Pro innerem Trainingsspender werden 10.000 Alive-Zellen ohne Zurücklegen
gezogen. Bei neun oder zehn Spendern sind es 90.000 bzw. 100.000
Trainingszeilen. Ein \code{StandardScaler} wird auf genau dieser
verketteten Stichprobe gefittet. ArcSinh wurde bereits auf Float32-
Markerwerten berechnet. Validierungszellen verwenden unveränderte
Trainingsmittel und -skalen. Alle drei C-Werte erhalten für denselben
Fold dieselbe Zellstichprobe und denselben Scaler.

Das Zellmodell lautet
\[
 g(z_i)=w^\top z_i+b.
\]
Positive Margen sprechen in der binären Klassenrichtung zunächst für
CMV+. Die Parameter sind Gewichte der \emph{standardisierten} Marker,
nicht direkt der Rohintensitäten. Ein einzelnes Gewicht lässt sich daher
nicht unabhängig von Transformation, Skala und anderen Markern biologisch
als Relevanzgröße lesen.

\code{LinearSVC} verwendet L2-Regularisierung und Squared-Hinge-Verlust.
Mit \(t_i=2y_i-1\in\{-1,+1\}\) und Intercept-Skalierung 1 entspricht
seine Optimierung hier der Form
\[
 \min_{w,b}\ \frac12(\|w\|_2^2+b^2)
     +C\sum_i\big[\max(0,1-t_i(w^\top z_i+b))\big]^2.
\]
Der Intercept wird in liblinear über eine synthetische konstante
Featurekomponente repräsentiert und bei diesen Defaults mitregularisiert.
Das unterscheidet ihn von den unbestraften CellCNN-Biases und dem
unbestraften glmnet-Intercept\cite{linearsvc}.

Für eine illustrative positive Zelle mit Marge 0,2 beträgt der
Squared-Hinge-Term \((1-0{,}2)^2=0{,}64\). Bei Marge 1,5 ist er null;
bei falscher Marge −0,5 beträgt er \((1+0{,}5)^2=2{,}25\).
\(C\) gewichtet diese Zellverluste gegenüber der Strafe.

Weitere Parameter sind \code{dual="auto"}, maximal 10.000 Iterationen
und Optimierungstoleranz \(10^{-4}\). Das Problem hat erheblich mehr Zellen
als Marker; die automatische Solverwahl nutzt dafür den passenden
Optimierungsweg. Der gesetzte Seed dient der Reproduzierbarkeit, ist aber
nicht mit einer biologischen Randomisierung gleichzusetzen.
\source{Notebook 04b, \code{prepare_inner_fold},
\code{select_svm_hyperparameters}; \code{LinearSVC}-Dokumentation\cite{linearsvc}.}

\clearpage
% Seite 26
\subsection*{C-Auswahl und die gelernte Spenderschwelle}
Getestet werden \(C\in\{0{,}01;0{,}1;1\}\). Je C wird eine neue SVM
in jedem der drei inneren Folds trainiert. Die Validierung bewertet alle
Zellen der jeweils ausgeschlossenen vier oder fünf Spender. Deren
Zellmargen werden zunächst je Spender aggregiert (nächste Seite), erst
danach wird ROC-AUC berechnet.

\[
 \overline{\mathrm{AUC}}(C)=\frac13\sum_{f=0}^2\mathrm{AUC}_f(C).
\]
Das ist das ungewichtete Mittel der drei Fold-AUCs, nicht eine einzige AUC
aller zusammengeführten OOF-Scores. Das größte Mittel gewinnt; bei
Gleichstand wird das kleinste C gewählt. Alle 900 exportierten
C-/Fold-/Split-Zeilen sind erhalten.

Im vorliegenden Lauf sind die AUCs der drei C-Werte in jedem inneren Fold
gleich. Deshalb wird in allen 100 äußeren Splits C=0,01 ausgewählt. Das
bedeutet \emph{nicht}, dass alle Modelle oder Zellmargen identisch sind.
AUC hängt nur von Rangvergleichen ab; kleine Wertänderungen ohne
Rangänderung verändern diese Metrik nicht. Mit vier oder fünf
Validierungsspendern ist ihre Auflösung zudem grob.

\textbf{Schwellenwahl aus inneren OOF-Scores.} Für das gewählte C liegt
für jeden der 14 äußeren Trainingsspender genau ein Score eines Modells
vor, das diesen Spender nicht trainiert hat. Auf diesen 14 Scores wird
\[
 J(\tau)=\operatorname{Sensitivität}(\tau)+
         \operatorname{Spezifität}(\tau)-1
        =\operatorname{TPR}(\tau)-\operatorname{FPR}(\tau)
\]
maximiert. Die Schwellen stammen aus \code{roc_curve}; unendliche
Schwellen werden ausgeschlossen. Bei Gleichstand gewinnt der erste endliche
Maximierer in der absteigenden Schwellenreihenfolge, also die höchste dort
angebotene Schwelle. Es werden keine äußeren Testlabels verwendet.

Erst danach wird eine neue finale SVM auf 10.000 Zellen aus jedem der
14 äußeren Trainingsspender trainiert. Ihr Scaler wird ebenfalls auf
dieser neuen Trainingsstichprobe gefittet. Die Schwelle bleibt der
zuvor bestimmte Wert. Trainingssampling nutzt Seed
\(s+20.000+100\); inneres Sampling \(s+20.000+f\).

\finding{Die Schwelle verbindet rohe Scores verschiedener innerer Modelle
und wird auf ein neu gefittetes Modell übertragen. Vergleichbare
Scoreskalen werden dabei angenommen; eine gesonderte Kalibrierung ist
nicht implementiert. Das kann Balanced Accuracy beeinflussen, nicht die
schwellenunabhängige ROC-AUC. Die historischen 14 OOF-Einzelscores je C
werden nicht exportiert; eine vollständige Neuberechnung der Schwellenwahl
würde innere Fits erfordern und ist hier nicht erfolgt.}
\source{\code{select_youden_threshold}, \code{select_svm_hyperparameters},
\code{fit_and_predict_outer_split}; \code{task4_svm_selection_gated_alive_full.csv}.}

\clearpage
% Seite 27
\subsection*{Top-1-\%-Aggregation und endgültige SVM-Vorhersage}
Für einen Testspender mit \(N_d\) Zellen berechnet die SVM alle Margen.
Die \(k_d=\max(1,\lceil0{,}01N_d\rceil)\) größten werden gemittelt:
\[
 S_d=\frac1{k_d}\sum_{i\in T_d}g(z_i),\qquad
 \widehat y_d=\mathbf1\{S_d\ge\tau\}.
\]
Hier wird \textbf{aufgerundet}, im Unterschied zu CellCNN. Bei 101 Zellen
würden zwei SVM-Margen gepoolt, CellCNN würde im entsprechenden Input nur
eine Antwort nehmen. Die realen Proben sind wesentlich größer; die
Rundungsregel wird dennoch exakt dokumentiert.

\textbf{Illustrative Rechnung.} Bei 1.000 Zellen werden zehn Margen
verwendet. Sind diese
\[1{,}1;1{,}2;1{,}2;1{,}3;1{,}3;1{,}4;1{,}4;1{,}5;1{,}6;2{,}0,\]
ist \(S_d=1{,}4\). Bei gelernter Schwelle 1,3 wird der Spender positiv
klassifiziert. Die anderen 990 Zellen beeinflussen diesen Mittelwert
nicht direkt. Sie können die Zugehörigkeit zur Top-Menge aber verändern,
wenn ihre Margen bei anderen Parametern steigen.

\textbf{Warum Top-Pooling?} Der vorab festgelegte Anteil soll seltene,
stark reagierende Zellen gegenüber der großen Hintergrundpopulation
hervorheben. Das ist eine Projektanpassung der Zellbaseline. Sie wurde
nicht anhand der äußeren Test-AUC optimiert. Ein alternatives Mittel über
alle Zellmargen ist kein Teil dieses Benchmarks.

Die \code{partition}-Operation findet die größten Werte, ohne alle Margen
vollständig zu sortieren. Bei gleichen Randwerten ist deren konkrete
Zellidentität für den Spenderscore irrelevant. Für Aufgabe 5 wird sie
hingegen relevant: Dort entscheidet zusätzlich der Originalereignisindex,
damit exakt dieselbe Zahl konkreter Zellen eindeutig gewählt werden kann.

\begin{table}[!ht]\centering\small
\caption{Gespeicherte SVM-Parameter und ihre Kontrollfunktion.}
\begin{tabular}{p{5cm}p{10cm}}\toprule
Bestandteil & Zweck\\\midrule
37 Markergewichte, Intercept & Zellmarge \(g(z)\) rekonstruieren\\
Scaler-Mittel und -Skalen & Dieselbe Trainingsstandardisierung anwenden\\
Klassenreihenfolge und Schwelle & Richtung und Spenderentscheidung prüfen\\
C, Seed, Gate, Zellzahl & Fit und Sampling dem richtigen Split zuordnen\\
Top-Anteil / Top-Zellzahl & Aggregationsregel und tatsächlich verwendete Anzahl prüfen\\
\bottomrule\end{tabular}
\end{table}

Alle 600 gespeicherten SVM-Spenderscores sind positiv. Dies widerspricht
einem funktionierenden Klassifikator nicht: Ein Mittel der höchsten Margen
liegt erwartbar deutlich höher als ein globales Zellmittel. Eine
nachträglich eingesetzte Schwelle null würde die ausgewählte
Spenderentscheidung verändern. Die gespeicherten Schwellen liegen
ungefähr zwischen 0,676 und 1,664.

\textbf{Interpretationsgrenze.} Das lineare Zellmodell lernt auf vielen
schwach beschrifteten Zellen; die getestete Vorhersage stammt von deren
Extrembereich. Ein gutes Spender-Ranking beweist deshalb weder ein korrektes
CMV-Label einzelner Zellen noch vollständige Wiederfindung eines seltenen
biologischen Subsets.
\source{\code{aggregate_top_fraction}, \code{fit_and_predict_outer_split};
\code{task4_svm_models_gated_alive_full.csv}, Vorhersagetabelle mit Top-Zellzahl.}

\clearpage
% Seite 28
\section{Klassifikationsergebnisse und ihre Aussagekraft}\label{sec:ergebnisse}
\subsection*{Vier Metriken beantworten verschiedene Fragen}
\textbf{ROC-AUC} bewertet Rangordnung. Mit \(n_+=2\) positiven und
\(n_-=4\) negativen Testspendern gilt
\[
 \mathrm{AUC}=\frac1{n_+n_-}\sum_{d:y_d=1}\sum_{e:y_e=0}
 \left[\mathbf1(S_d>S_e)+\tfrac12\mathbf1(S_d=S_e)\right].
\]
Es gibt nur acht positiv-negative Paare. Ohne Gleichstände springt AUC in
Schritten von 0,125, mit halben Paarpunkten auch von 0,0625. Eine monotone
Scoretransformation ändert die Rangfolge nicht.

\textbf{Average Precision (AP)} gewichtet Precision mit Recall-Zuwächsen:
\(\mathrm{AP}=\sum_j(R_j-R_{j-1})P_j\). Gleiche Scores werden gemeinsam
als Schwelle behandelt. \textbf{Trapezoidale PR-AUC} integriert dagegen
linear zwischen den Kurvenpunkten:
\(\sum_j(R_j-R_{j-1})(P_j+P_{j-1})/2\).
Beide Größen liegen zwischen null und eins, sind aber nicht gleich.
\textbf{Balanced Accuracy} ist
\(\tfrac12[TP/(TP+FN)+TN/(TN+FP)]\) an der jeweiligen Modellschwelle.

\tab{metriken}{Vollständige Metrikübersicht über 100 gemeinsame äußere
Splits. Q1/Q3 sind Quartile, SD die Stichprobenstandardabweichung über
Splits. Die 1.200 Einzelwerte und ihre Zusammenfassungen wurden für den
Bericht unabhängig nachgerechnet.}{tab:metrics}{}

\textbf{Konstantes Modell als Handrechnung.} Sechs gleiche Scores ordnen
kein Paar: AUC=0,5. Bei zwei Positiven ist AP=\(2/6=1/3\).
Die trapezoidale PR-Kurve verbindet den Startpunkt \((R=0,P=1)\) mit
\((R=1,P=1/3)\), daher ist ihre Fläche \((1+1/3)/2=2/3\).
Die gleiche Zahl kann somit eine überraschend hohe PR-AUC ergeben, obwohl
keine Ranginformation vorliegt. Dieser Unterschied begründet, AP zusätzlich
zu berichten. Der positive Anteil ist keine passende Referenzlinie für
diese trapezoidale Fläche.
\source{Notebook 04e; unabhängige Umsetzung \code{independent_metrics}
in \code{src/detail_report_assets.py}; tests mit Ranggleichständen und
konstanten Scores.}

\clearpage
% Seite 29
\subsection*{Leistung auf denselben Splits vergleichen}
\fig{roc_auc}{.88}{ROC-AUC auf Spenderebene für den Full-Hauptlauf.
Boxen: Quartile, Mittellinie: Median; überlagerte Punkte: einzelne Splits.
Die Linie bei 0,5 bezeichnet fehlende Rangtrennung. Punkte können wegen der
groben Metrikauflösung exakt übereinanderliegen.}{fig:auc}
\fig{auc_differenzen}{.88}{Gepaarte Differenzen: pro identischem Split wird
die AUC der Baseline von der CellCNN-AUC abgezogen. Positive Werte sprechen
in diesem Split für CellCNN. Null ist gleiche AUC. Die Paarung berücksichtigt,
dass alle Methoden dieselben sechs Testspender beurteilen.}{fig:diff}
\tab{differenzen}{Paarvergleich einschließlich der Zahl positiver, gleicher
und negativer Splitdifferenzen. Dies sind beschreibende Häufigkeiten,
keine 100 unabhängigen experimentellen Wiederholungen.}{tab:diff}{}

CellCNN liegt in 73 Splits vor Citrus, in neun gleich und in 18 dahinter.
Gegenüber SVM sind es 41, 35 und 24. Die mediane gepaarte Differenz zu
Citrus beträgt 0,25, zu SVM null. CellCNN hat eine höhere mediane Balanced
Accuracy als SVM, aber diese verwendet unterschiedliche Entscheidungsregeln
und hängt von Schwellen ab.

\finding{Belegt ist eine unter den gewählten Bedingungen stärkere
Rangtrennung von CellCNN gegenüber Citrus im Median. Gegenüber SVM wird
kein genereller Rangvorteil nachgewiesen. Quartile sind keine
Konfidenzintervalle; ein ungepaarter oder auf unabhängigen Splits beruhender
Signifikanztest wird nicht nachträglich behauptet.}

\clearpage
% Seite 30
\subsection*{Batchkontrolle, Laufzeit und allgemeine Grenzen}
Messtag und Instrumentversion können Expressionsmuster beeinflussen.
Die QC-Heatmap zeigt mögliche blockartige Verschiebungen. Als beschreibende
Kontrolle wird deshalb pro äußerem Trainingssplit die CMV+-Häufigkeit jedes
Messtags bzw. Instruments berechnet. Ein Testspender erhält die geschätzte
Rate seiner Gruppe; bei unbekannter Gruppe die Gesamttrainingsrate.
Dieser Kontrollscore ist keine vierte Hauptmethode.

\begin{table}[!ht]\centering
\caption{CMV-Labelzahlen je Messtag und Metadaten-Kontrollleistung.
Die Gruppenverteilung betrifft 20 unabhängige Spender.}
\begin{tabular}{lrr}\toprule
Messtag & CMV− & CMV+\\\midrule
1. Mai 2012 & 3 & 3\\
8. Mai 2012 & 3 & 2\\
12. Juni 2012 & 5 & 4\\\bottomrule
\end{tabular}\qquad
\begin{tabular}{lrr}\toprule
Kontrollvariable & AUC-Median & AUC-Mittel\\\midrule
Messtag & 0,250 & 0,306\\
Instrument & 0,375 & 0,376\\\bottomrule
\end{tabular}
\end{table}

Beide CMV-Gruppen kommen an jedem Messtag vor. Die niedrigen Kontroll-AUCs
sind kein Beweis vertauschter Labels: Wenn viele positive Spender eines
Batches im Training landen, bleiben in dessen komplementärem Test eher
negative übrig. In der nahezu balancierten kleinen Kohorte kann die
Trainingsrate dann gegenläufig zum Testlabel rangieren.

Frühere Prüfnotizen untersuchen diesen Mechanismus auch über alle möglichen
Testmengen und permutierte Labels. Diese Zusatzversuche wurden für den
Bericht nicht erneut ausgeführt; ihre historischen Detailzahlen werden
nicht als neue Prüfergebnisse ausgegeben. Die hier gezeigte Tabelle
beschreibt die bestehenden 100 Splits und die Originalmetadaten.

\textbf{Zulässige Schlussfolgerung:} Die Gruppen sind über Messtage
annähernd gleich verteilt; ein reiner Metadatenscore zeigt in diesem
Splitprotokoll keine positive Vorhersageleistung. \textbf{Nicht zulässig:}
„Batch-Effekte sind ausgeschlossen.“ Instrument und Messtag sind außerdem
gekoppelt und keine unabhängigen Kontrollen. Eine Validierung auf einem
neuen Messtag oder einer externen Kohorte liegt nicht vor.

\textbf{Keine faire Laufzeit-Rangliste.} CellCNN lief auf einer GPU,
SVM und Citrus auf CPUs; konkurrierende Prozesse und verschiedene
Implementierungen beeinflussen Zeiten. Verfügbare Protokollzeiten sind
keine kontrollierten Messungen vergleichbarer Rechenbudgets. Deshalb
wird kein vermeintlich exakter Geschwindigkeitssieger berichtet.

\textbf{Weitere Grenzen des Benchmarks.} Die Testmenge umfasst immer nur
sechs Spender. Die drei Verfahren verwenden verschiedene Zellzahlen,
Suchräume, Auswahlmetriken und finale Trainingsgrößen. Der Benchmark
vergleicht diese dokumentierten Pipelines, nicht einen isolierten
Architekturunterschied unter identischen Optimierungsbedingungen.
Es wurden weder eine vollständige Gate-Sensitivitätsanalyse noch
Markerablationen oder eine Moment-Baseline durchgeführt.
\source{Notebook 04e, Abschnitt 7; \code{task4_batch_control.csv},
\code{task4_batch_label_counts.csv}; historische Notiz
\code{untersuchung/AUFGABE_4_PRUEFUNG.md}.}

\clearpage
% Seite 31
\section{Aufgabe 5: welche Zellen unterstützen die Modelle?}\label{sec:task5}
\subsection*{Identität, unveränderte Modelle und zulässige Testauftritte}
Aufgabe 5 soll ein quantitatives Zellsubset bestimmen und auf einer
Dimensionsreduktionskarte zeigen. Dafür werden keine neuen Klassifikatoren
trainiert. Wir laden die gespeicherten SVM-/CellCNN-Parameter,
Citrus-Koeffizienten und die t-SNE-Karte mit Perplexität 30 aus Aufgabe 2.
Die Karte zeigt dieselben 500 Zellen je Spender, insgesamt 10.000.

Die Verbindung erfolgt über Zell-ID, Spender-ID und nullbasierten
Originalereignisindex. Das Laden kontrolliert das ursprüngliche Sampling,
die Reihenfolge und eine gemeinsame SHA-256-Prüfsumme über IDs, Marker
und rohe Kartenwerte. Ein stilles Zusammenfügen bloß gleich langer
Tabellen wäre unzureichend: Es könnte Scores anderen Zellen zuordnen.
Die Kartenidentität wurde beim Bericht erneut geprüft.

Für eine Zelle \(i\) aus Spender \(d\) sei
\[
 \mathcal S_d=\{s:\ d\text{ ist äußerer Testspender in Split }s\}.
\]
Nur diese Modelle dürfen die Zelle für die berichtete Selektion bewerten.
Daten aus \emph{anderen} äußeren Splits dürfen denselben Spender beim
Training enthalten; sie zählen für dessen OOF-Interpretation dann nicht.
Die t-SNE-Geometrie wurde explorativ auf allen Spendern angepasst und ist
selbst kein OOF-Prädiktor. Ihre Verwendung zur Darstellung ändert die
getrennten Zellbewertungen nicht, bietet aber auch keine unabhängige
Validierung ihrer räumlichen Struktur.

Für binäre Auswahlindikatoren \(I_{is}^{+}\) und \(I_{is}^{-}\) werden
\[
 f_i^+=\frac{\sum_{s\in\mathcal S_d}I_{is}^+}{|\mathcal S_d|},\qquad
 f_i^-=\frac{\sum_{s\in\mathcal S_d}I_{is}^-}{|\mathcal S_d|}
\]
berechnet. Die Nenner variieren hier zwischen 12 und 42 Testauftritten
je Spender. Alle zulässigen Modelle zählen, auch Nullmodelle. Fehlende
Modellbewertungen sind dagegen Fehler und dürfen nicht stillschweigend
als Nullselektion gelten.

\textbf{Beispiel.} Wird eine Zelle in drei von zwölf zulässigen Modellen
positiv gewählt, ist \(f_i^+=0{,}25\). Drei von 42 entsprechen nur
0,0714. Eine gleiche Anzahl positiver Wahlen ist deshalb noch keine gleiche
Stabilität. Auch gleiche \emph{Häufigkeit} verschiedener Methoden bedeutet
keine gleiche Effektstärke, denn deren Auswahlregeln unterscheiden sich.

Am Ende gibt es 30.000 Zeilen: eine pro Methode und Karten-Zelle. Sie
enthalten Auswahlzählungen, Nenner, Häufigkeiten und vorhandene
Koordinaten. Das ist eine Zusammenfassung vieler Modelle und kein
zusätzlich gelerntes Ensemble mit neu gemessener Klassifikationsleistung.
\source{\code{load_inputs}, \code{align_projection}, \code{validate_models},
\code{aggregate_frequencies} in \code{src/task5_interpretation.py}.}

\clearpage
% Seite 32
\subsection*{CellCNN: Phänotypregel ist nicht Poolingmitgliedschaft}
Das gespeicherte Netz berechnet \(r_{ik}=\max(0,W_kz_i+b_k)\) wie beim
Training. Für jeden ausgewählten Filter wird zunächst sein Maximum
\(M_{sk}\) über \emph{alle ursprünglichen Zellen seiner tatsächlichen
inneren Trainingsspender} bestimmt. Es werden also nicht nur die 3.000
Zellziehungen einzelner Inputs und auch nicht die äußeren Testspender
als Referenz verwendet.

Die papernahe Halbmaximumregel lautet
\[
 H_{isk}=\mathbf1\{M_{sk}>0\}\,
          \mathbf1\{r_{isk}>M_{sk}/2\}.
\]
Bei genau halbem Maximum ist eine Zelle nicht ausgewählt. Ein Filter,
der auf Trainingszellen nur null antwortet, definiert kein Subset.
Mit \(\Delta V_{sk}=V_{1k}-V_{0k}\) folgt
\[
 I_{is}^+=\mathbf1\{\exists k:\ H_{isk}=1,\ \Delta V_{sk}>0\},
 \quad
 I_{is}^-=\mathbf1\{\exists k:\ H_{isk}=1,\ \Delta V_{sk}<0\}.
\]
Für den Kontrast wird hier keine zusätzliche numerische Nulltoleranz
verwendet; sie ist bei Citrus eine andere explizite Regel.

\textbf{Illustration.} Ein Filter mit Trainingsmaximum 8 und positivem
Kontrast wählt eine Testzelle mit Antwort 5 positiv aus. Ein zweiter Filter
mit Maximum 6, Antwort 4 und negativem Kontrast wählt dieselbe Zelle
negativ. Beide Indikatoren dürfen gleichzeitig eins sein. Im gespeicherten
Lauf betrifft eine solche gleichzeitige Auswahl mindestens einmal sieben
Karten-Zellen. Über \emph{verschiedene} Splits können Richtungen auch bei
viel mehr Zellen wechseln.

Die Halbmaximumregel beschreibt einen zu einem Filter passenden Phänotyp.
Sie besagt nicht, dass die Zelle in einem tatsächlich für die
Spendervorhersage gezogenen Testinput lag. \code{exposed_models} zählt,
in wie vielen zulässigen Modellen die Karten-Zelle wenigstens in einem
der fünf Testinputs vorkam. Das ist wiederum nicht identisch mit der
Top-k-Mitgliedschaft: Auch eine vorkommende Zelle kann unter der
Poolinggrenze liegen. Die Pipeline exportiert keine vollständigen
zellweisen Top-k-Beiträge aller fünf Inputs.

Der tatsächliche Inputlogit erfüllt hingegen die prüfbare Identität
\[
 \ell_1-\ell_0=(d_1-d_0)+\sum_k\Delta V_k\,
       \frac1q\sum_{i\in T_k}r_{ik}.
\]
Hier taucht nur die pro Input und Filter gepoolte Menge \(T_k\) auf.
Aufgabe 5 kontrolliert diese Identität, berichtet auf der Karte aber die
Halbmaximumselektion. Zudem ist die Spendervorhersage der Mittelwert von
fünf \emph{Wahrscheinlichkeiten}; sie besitzt nicht einfach den Logit
eines einzigen gepoolten Gesamtinputs.

\finding{Die Aussage „jede markierte Zelle hat die konkrete
Spendervorhersage mitbestimmt“ ist für CellCNN nicht belegt. Belegt ist
ihre wiederholte Übereinstimmung mit gerichteten Filtern von Modellen,
die ihren Spender nicht trainiert haben.}
\source{\code{cellcnn_masks}, \code{SavedCellCNN}, \code{run_python_methods};
\code{task5_filter_thresholds.csv}; Paper-Halbmaximumregel\cite{cellcnn}.}

\clearpage
% Seite 33
\subsection*{Citrus: addierte Clusterbeiträge und Baumrekonstruktion}
Citrus wurde nicht mit vollständigen Baumobjekten exportiert. Für die
Interpretation werden deshalb die finalen Trainingsbäume aus denselben
1.000 Zellen je ursprünglichem Trainingsspender und denselben R-Seeds
rekonstruiert. Das ist eine erneute \emph{unüberwachte Baumkonstruktion}
als Teil des bereits ausgeführten Aufgabe-5-Laufs, kein neues
Regressions- oder CV-Training. Dieser Bericht selbst führt auch diese
Baumrekonstruktionen nicht erneut aus.

Das R-Skript prüft Originaldaten, Splitzuordnung, R- und Paketversionen,
Quellcommits, ausgewählte Cluster-IDs, Mindestgrößen und die exportierten
Markerzentroiden. Von 100 Modellen besitzen 85 wirksame Cluster und
benötigen einen rekonstruierten Baum. Bei den 15 Nullmodellen sind alle
Zellbeiträge null; dort ist ein Baum für die Auswertung nicht erforderlich.

Für eine Karten-Zelle wird die nächste Trainingszelle mit dem nativen
Citrus-Mapping ermittelt. Der Zellscore ist
\[
 a_{is}=\sum_{k:\,|\beta_{sk}|>10^{-10}}
         \beta_{sk}\,\mathbf1\{i\text{ wird in }C_{sk}\text{ gemappt}\}.
\]
Alle überlappenden Mitgliedschaften zählen. Die positive Richtung gilt
bei \(a_{is}>10^{-10}\), die negative bei \(a_{is}<-10^{-10}\).
Eine nahezu vollständige Aufhebung positiver und negativer Gewichte
kann deshalb trotz Clusterzugehörigkeit zu keiner gerichteten Auswahl führen.

Für die 1.000 tatsächlich zur Testvorhersage gezogenen Zellen eines
Spenders gilt die lineare Identität
\[
 \frac1{1000}\sum_i a_{is}=\sum_k\beta_{sk}F_{dk}.
\]
Das erklärt den Zusammenhang zwischen Zell- und Spenderrepräsentation.
Die Kartenstichprobe umfasst jedoch 500 andere, fest gespeicherte Zellen
pro Spender; ihr Mittelwert ist nicht automatisch derselbe Testfeaturescore.
\code{exposed=True} bedeutet hier, dass das Phänotypmapping für alle
Karten-Zellen definiert ist. Es bestätigt nicht deren Zugehörigkeit zur
ursprünglichen 1.000-Zell-Teststichprobe.

\textbf{Was lässt sich ohne Intercept prüfen?} Das gespeicherte
\(\beta_0\) fehlt. Aufgabe 5 rekonstruiert daher keine absoluten
Citrus-Wahrscheinlichkeiten. Mit \(L_d=\operatorname{logit}(p_d)\)
und \(A_d=\sum_k\beta_kF_{dk}\) prüft es
\[
 L_d-\overline L=A_d-\overline A,
 \qquad L_d-L_e=A_d-A_e.
\]
Der unbekannte gemeinsame Intercept fällt weg. Zusätzlich wird die
Feature-Summe unabhängig gegen die gemittelten Zellbeiträge kontrolliert.
Das Feld \code{decision_matches} prüft für Citrus nur, ob die gespeicherte
Klasse zur gespeicherten Wahrscheinlichkeit und Schwelle 0,5 passt.
Es ist keine unabhängig rekonstruierte absolute Entscheidung.
\source{\code{src/task5_citrus.R}: \code{cluster_cell_scores},
\code{centered_logit_errors}, \code{run_task5_citrus};
\code{task5_citrus_cluster_checks.csv}.}

\clearpage
% Seite 34
\subsection*{SVM: exakte Poolingzellen und Richtung an der Spenderschwelle}
Aus den gespeicherten Parametern werden alle Zellmargen eines Testspenders
rekonstruiert. \code{exact_top_indices} wählt genau
\(\lceil0{,}01N_d\rceil\) der höchsten Margen. Bei Gleichstand am
Rand entscheiden die kleineren Originalereignisindizes. Diese zusätzliche
Regel ändert nicht den aggregierten Score, sorgt aber für eindeutige
Zellidentitäten.

Nur innerhalb dieser Menge \(T_d\) wird die Richtung relativ zur
gelernten Spenderschwelle \(\tau_s\) bestimmt:
\[
 I_{is}^+=\mathbf1\{i\in T_d\}\mathbf1\{g_s(z_i)>\tau_s\},\qquad
 I_{is}^-=\mathbf1\{i\in T_d\}\mathbf1\{g_s(z_i)<\tau_s\}.
\]
Eine exakt auf der Schwelle liegende gepoolte Zelle zählt zu keiner
Richtung. Pro Modell sind positiv und negativ disjunkt. Über mehrere
Modelle kann eine Zelle trotzdem in beiden Richtungen auftreten.

Der Zusammenhang zur Entscheidung ist unmittelbar:
\[
 S_d-\tau_s=\frac1{k_d}\sum_{i\in T_d}(g_s(z_i)-\tau_s).
\]
Positive und negative Beiträge innerhalb der Top-Menge ergeben gemeinsam
die Lage des Spenders zur Schwelle. Eine als negativ markierte SVM-Zelle
ist also nicht die Zelle mit der niedrigsten Marge der gesamten Probe.
Sie gehört weiterhin zu den höchsten 1\%, liegt aber unter der
Spenderschwelle. Der Begriff „negativ“ ist ohne diese Definition leicht
missverständlich.

\tab{selektion}{Auswahl auf den 10.000 Karten-Zellen. „Zellen“ zählt pro
Methode und Richtung die mindestens einmal gewählten Zellen. „Mittl. Anteil“
ist dagegen die mittlere Auswahlhäufigkeit innerhalb eines Spenders,
anschließend über Spender gemittelt. Min.–Max. zählt mindestens einmal
gewählte Karten-Zellen pro Spender.}{tab:selektion}{}

\textbf{Warum ist der SVM-Anteil nicht genau 1\%?} Die Top-1-\%-Regel gilt
für alle Zellen eines Testspenders. Die Tabelle zeigt nur die 500 zufällig
auf die Karte gelangten Zellen je Spender. Außerdem werden die Richtungen
getrennt und über Modelle gemittelt. „Mindestens einmal gewählt“ und
„durchschnittlich pro Modell gewählt“ sind verschiedene Größen.

Bei positiver SVM-Selektion tragen alle 20 Spender Profile bei, teilweise
aber nur \emph{eine einzige Karten-Zelle}. Eine starke Markeranreicherung
kann deshalb ein sehr kleines, fragiles Teilbild beschreiben. Sie belegt
nicht, dass die SVM die biologisch korrekte Population besser wiederfindet.
\source{\code{exact_top_indices}, \code{svm_cell_selection},
\code{run_python_methods}, \code{summarize_profiles};
\code{task5_selection_summary.csv}.}

\clearpage
% Seite 35
\subsection*{Positive und negative Auswahl auf identischer t-SNE-Karte}
\fig{selektion_positive}{1}{Positive Auswahlhäufigkeit \(f_i^+\) auf der
unveränderten t-SNE-30-Karte. Je Methode dieselben 10.000 Alive-Zellen,
500 je Spender. Grau: nie positiv gewählt; farbig: Häufigkeit über alle
äußeren Testmodelle des Spenders, gemeinsame Skala 0--1. Methodenregeln
siehe vorherige Seiten; gleicher Farbwert bedeutet keine gleiche Wirkung.}{fig:pos}
\fig{selektion_negative}{1}{Negative Auswahlhäufigkeit \(f_i^-\), identische
Koordinaten und Skala. Bei SVM ist dies die unterhalb der Spenderschwelle
liegende Teilmenge der höchsten 1\% Zellmargen. Bei CellCNN können sich
positive und negative Filterselektionen überlagern; Citrus verwendet die
Nettosumme der Clustergewichte.}{fig:neg}
\fig{abdeckung}{.86}{Links: Zahl äußerer Testauftritte pro Spender, also der
Nenner der Zellhäufigkeiten. Rechts: kumulative Verteilung positiver
Häufigkeiten unter den mindestens einmal positiv gewählten Karten-Zellen.
Nicht gewählte Zellen sind rechts ausgeschlossen; die Kurven haben daher
unterschiedlich große Bezugsgruppen.}{fig:abdeckung}

CellCNN wählt 1.313, Citrus 7.678 und SVM 219 Karten-Zellen mindestens
einmal positiv. Breite Citrus-Abdeckung und schmale SVM-Abdeckung folgen
teilweise bereits aus den Auswahlregeln. Die Häufigkeiten dürfen deshalb
nicht als ein gemeinsamer Zellklassifikationsscore gegeneinander in
ROC-Kurven gewertet werden: Wahrheitslabels für die gesuchten Zellsubsets
fehlen. Der OOF-Bezug begrenzt Überanpassung bei der Zellbewertung,
erzeugt aber keine biologische Ground Truth.
\source{\code{task5_cell_scores.csv}, \code{task2_embeddings.csv};
keine erneute t-SNE-Anpassung.}

\clearpage
% Seite 36
\subsection*{Markerprofile: wie die Zahlen entstehen und was sie bedeuten}
Für Methode, Richtung, Spender und Marker werden die ArcSinh-Werte der
500 Karten-Zellen mit deren Auswahlhäufigkeit gewichtet. Nach Sortierung
wird der kleinste Wert gewählt, bei dem die kumulierte positive
Gewichtssumme mindestens die Hälfte der gesamten Gewichtssumme erreicht.
Das ist ein gewichteter Median mit unterem Wert bei exaktem Halbgewicht.
Ohne positives Gewicht entsteht ein fehlendes Profil, keine Null-Expression.
Danach wird der Median der vorhandenen Spenderprofile berechnet.

Die Kartenreferenz verwendet ungewichtete Mediane je Spender, dann den
Median über Spender. Alle Profile verwenden Float64-ArcSinh ohne Modell-z;
Vorhersagen behalten ihre ursprüngliche Float32-/Trainingsscaler-Arithmetik.
\fig{selektionsprofile}{.82}{Positive und negative Profile gegenüber der
Kartenreferenz. Häufigkeitsgewichtete Mediane innerhalb der Spender,
anschließend Median über unterstützte Spender, in ArcSinh-Einheiten.
Dies sind marginale Markerzusammenfassungen; benachbarte hohe Werte
beweisen keine gemeinsame Expression in derselben Zelle.}{fig:profile5}
\tab{markerwerte}{Ausgewählte Marker der positiven Profile als genaue
Zahlen. Alle 37 Marker und beide Richtungen sind in
\code{task5_marker_profiles.csv} dokumentiert.}{tab:marker5}{}

SVM zeigt gegenüber der Kartenreferenz erhöhte NKG2C-, CD57- und
CD56-Werte. CellCNN zeigt diese Richtung ebenfalls, für NKG2C/CD57
schwächer. Beide Profile behalten jedoch hohes CD3. Das unterstützt
NK-assoziierte Rezeptormuster mit T-kompatiblen oder heterogenen Anteilen,
\emph{keinen isolierten CD3-negativen NK-Zelltyp}. Citrus ist im positiven
Profil CD3/CD8-reich, mit niedrigerem CD56/NKG2C, und beschreibt keine
schmale stabil wiederkehrende NK-Population. Negative Selektionen liefern
zusätzliche Gegenprofile, aber ebenfalls keine eindeutigen Zelltyp-Labels.

\finding{Die Aufgabe ist als quantitative modellbezogene Selektion mit
Visualisierung bearbeitet. Ein biologisch validiertes Wiederfinden genau
der NKG2C+/CD57+-NK-Population des Papers oder dessen Zentroidgruppen
ist damit nicht nachgewiesen.}
\source{Aufgabe 5: \code{weighted_median}, \code{summarize_profiles}; Profile jetzt nachgerechnet.}

\clearpage
% Seite 37
\section{Prüfnachweise, Vollständigkeit und Reproduzierbarkeit}\label{sec:pruefung}
\subsection*{Was jetzt geprüft wurde und was aus dem Analyseprotokoll stammt}
Der deutsche Export führt keine Modellfits und keine neue Dimensionsreduktion
aus. Er liest FCS-Dateien, rekonstruierte Stichproben und bestehende Tabellen,
prüft deren Zuordnung und berechnet Berichtszahlen neu. Folgende Ebenen
sind auseinanderzuhalten:
\begin{table}[!ht]\centering\small
\caption{Nachweisstatus dieses Detailberichts. „Jetzt“ bedeutet bei seiner
Erstellung ausgeführt; ein vorhandener Nachweis ist kein neuer Trainingslauf.}
\begin{tabular}{p{4.5cm}p{2.6cm}p{7.8cm}}\toprule
Gegenstand & Status & Umfang und Grenze\\\midrule
FCS-Inventar und Werte & Jetzt geprüft & 40 Dateien, Endlichkeit der 37 Marker, Eventzahlen; Originaldateien nur gelesen\\
Kartenidentität & Jetzt geprüft & 10.000 Eventindizes und Rohwertprüfsumme; Koordinatenzuordnung aller drei Methoden\\
Spendersplits / Labels & Jetzt geprüft & 2.000 Zuordnungen und 600 Testvorhersagen je Methode\\
Metriken & Jetzt geprüft & 1.200 Werte über Rangpaare, Schwellen und Klassenquoten; Zusammenfassungen abgeglichen\\
Auswahlhäufigkeiten & Jetzt geprüft & 30.000 Zellzeilen, spenderabhängige Nenner, Richtungszählungen und Zusammenfassungen\\
Markerprofile & Jetzt geprüft & 5.180 gewichtete bzw. Referenz-Spenderwerte aus Original-Kartenzellen nachgerechnet\\
Aufgabe-5-Provenienz & Jetzt geprüft & 48 SHA-256-Nachweise zu Eingaben, Artefakten, Umsetzung und Ausgaben\\
Netz-/SVM-Inferenz & Vorhandener Nachweis & Je 600 Rekonstruktionen; Artefaktidentität geprüft, nicht erneut vollständig inferiert\\
Citrus-Bäume / Logits & Vorhandener Nachweis & 85 Bäume; 15 Nullmodelle; 600 relative Logitprüfungen, kein exportierter Intercept\\
Biologische Richtigkeit & Offen & Keine Zelltyp-Ground-Truth oder externe Validierung\\
\bottomrule\end{tabular}
\end{table}
\tab{rekonstruktion}{Maximale Fehler des gespeicherten Aufgabe-5-Prüflaufs,
nicht eines neuen Benchmarktrainings. Bei CellCNN betrifft der Rechenfehler
die Logitidentität, bei SVM die Schwellenzentrierung, bei Citrus zentrierte
Logits. Diese Größen haben unterschiedliche Bedeutungen.}{tab:checks}{}

Die Prüftoleranzen sind SVM-Score \(10^{-10}\), CellCNN-Wahrscheinlichkeit
\(10^{-6}\), CellCNN-Logitidentität \(10^{-5}\), Citrus-relative Logits
\(10^{-8}\). Die maximale paarweise Citrus-Logitabweichung ist etwa
\(3{,}0\times10^{-12}\); die Identität von Häufigkeits- und Zellbeiträgen
weicht höchstens etwa \(1{,}8\times10^{-15}\) ab. Kleine Rechenfehler
belegen numerische Konsistenz, nicht richtige biologische Annotation.

\textbf{Historische Befunde einordnen.} Die Notiz vom 5. September enthält
auch Probleme des Ausgangsstands: fehlende Konfigurationsnachweise,
unvollständige Modellparameter und numerische Nullclusterzählungen.
Die aktuellen Artefakte enthalten die entsprechenden Ergänzungen.
Diese früheren Probleme werden nicht als weiterhin bestehende Fehler
übernommen. Die hier neu benannte innere Lambda-Abhängigkeit und die
Interpretationsgrenzen bleiben ausdrücklich sichtbar.
\source{Beim Bericht bestanden alle 29 Python-Tests und beide R-Testskripte.
Drei neue Handrechentests prüfen Metriken einschließlich Ranggleichständen;
die übrigen Tests prüfen bestehende Analyse- und Exportfunktionen.}

\clearpage
% Seite 38
\subsection*{Anforderungen, offene Punkte und reproduzierbarer Arbeitsstand}
\begin{table}[!ht]\centering\small
\caption{Abgleich mit sämtlichen Aufgaben des bereitgestellten Projektblatts.}
\begin{tabular}{p{2cm}p{6.7cm}p{6.2cm}}\toprule
Anforderung & Umsetzung und Beleg & Grenze / offener Punkt\\\midrule
1: Literatur & Überblick zu Paper, CellCNN und Baselines, S. 3 & Kein Anspruch auf vollständige unabhängige Literaturübersicht\\
2: Karten & PCA/t-SNE/UMAP, acht Parameterbeispiele; Abb.\ \ref{fig:projektionen} & Kleines Raster, eine Stichprobe und ein Seed\\
2: Vergleich & Alle 28 Kartenpaare und 37D-Bezug; Abb.\ \ref{fig:jaccard} & Lokale Geometrie, kein Zelltypbeweis\\
3: Clustering & K-Means, Ward, Leiden, neun Lösungen; Tabelle\ \ref{tab:cluster} & Biologisch bester Ansatz nicht eindeutig validierbar\\
3: Zelltypen & Markerprofile und Annotation aller ausgewählten Cluster & Plausibilität; Referenzlabels fehlen\\
4: Methoden & Drei dokumentierte Pipelines, gleiche äußere Splits & Unterschiede in Sampling, Auswahlziel und finaler Trainingsmenge\\
4: Bewertung & Vier Metriken, gepaarte AUC-Differenzen & 20 Spender; keine externe Kohorte oder faire Laufzeitmessung\\
5: Zellscore & Gerichtete methodenspezifische Selektion und OOF-Häufigkeit & CellCNN-Halbmaximum nicht identisch mit tatsächlichem Poolingbeitrag\\
5: Visualisierung & Gleiche t-SNE-30-Karte, beide Richtungen, Markerprofile & Keine bewiesene Wiedergewinnung des Paper-NK-Subsets\\
6: Bonus & Als offen dokumentiert & Keine Erweiterung trainiert oder bewertet\\
\bottomrule\end{tabular}
\end{table}

\textbf{Wie der Bericht reproduziert wird.} Zuerst müssen die bestehenden
Aufgaben-2--5-Ergebnisdateien und Originaldaten lokal verfügbar sein. Der
separate Export prüft Eingaben und erzeugt ausschließlich
\code{report/detail_assets}. Der normale LaTeX-Build benötigt diese
fertigen Assets, keine Python-/R-Analyse und keinen Datensatz.
\begin{quote}\small\ttfamily
python -m src.detail\_report\_assets\par
latexmk -cd -pdf -outdir=build report/detailbericht\_de.tex\par
latexmk -cd -pdf -outdir=build report/main.tex
\end{quote}
Im vorhandenen System ist Python 3.12 über \code{.venv/bin/python}
verfügbar. Methodenversionen stehen in den gespeicherten Konfigurationen;
der Full-Lauf verwendete unter anderem scikit-learn 1.9.0 und
PyTorch 2.14.0 mit CUDA. Die R-Umgebung und festgelegte Citrus-Quellstände
stehen in der Projekt-README. Der PDF-Build benötigt zusätzlich deutsche
Babel-/Trennmuster. \code{quellen.json} dokumentiert die gelesenen
Tabellenprüfsummen, \code{pruefung.json} den Umfang des Berichtsexports.

\textbf{Was für eine stärkere Schlussfolgerung fehlen würde.} Eine externe
Spenderkohorte, Referenzzelltypen und gemeinsame Expressionsprüfungen der
selektierten Zellen wären für biologische Bestätigung nötig. Strengere
innere Citrus-Abstimmung, SVM-Kalibrierung und CellCNN-Ablationen wären
separate methodische Experimente. Sie sind keine versteckten Bestandteile
des aktuellen Berichts und werden nicht als bereits erledigt dargestellt.

Die vorliegenden Ergebnisse tragen eine vorsichtige Schlussfolgerung:
Die Pipelines erkennen in dieser Kohorte CMV-bezogene Ranginformation;
ihr Verhältnis zu einem eng definierten biologischen Zellsubset bleibt
teilweise offen. Genau diese Trennung macht die Arbeit überprüfbar.

\clearpage
% Seite 39
\section*{Literatur und technische Primärquellen}
\addcontentsline{toc}{section}{Literatur und technische Primärquellen}
\renewcommand{\bibsection}{}
\renewcommand{\bibfont}{\small}
\setlength{\bibsep}{3pt}
\bibliographystyle{unsrtnat}
\bibliography{references,detail_references}
\vspace{3mm}
\subsection*{Lokale Quellen und Zuordnung}
Das Aufgabenblatt \code{Group_projects_ssbi_2026.pdf}, Projekt 1,
legt Anforderungen 1--6 fest. Die lokale \code{CellCNN.pdf} enthält den
Originalartikel. Beide werden unverändert behandelt. Berichtsgrundlage ist
der versionierte Analysecode; die durch Analyseläufe erzeugten Original-
CSV-Dateien bleiben gemäß \code{.gitignore} lokal.

Für die Methodendetails sind die Abschnitte und Funktionsnamen an den
betreffenden Stellen angegeben. Das verhindert eine Abhängigkeit von
verschiebbaren Notebook-Ausgabezeilen. Die Tabellen des Detailberichts
werden aus den genannten CSVs erzeugt und sind keine manuell geschätzten
Ablesewerte aus Plots. Kleinere Rundungen dienen nur der Lesbarkeit.

Die beiden offiziellen CellCNN-Notebooks wurden am 6. September 2026
lesend vom Master-Branch abgerufen. Dies ist kein unveränderlich gepinnter
Bibliotheksstand; für ihre hier genannten expliziten Beispielargumente
wurden die Notebookquellen selbst kontrolliert. Die tatsächlich
verwendete Python-Neuimplementierung ist separat im Projekt versioniert.
Die Citrus- und Rclusterpp-Implementierungen sind hingegen auf die im
Methodenteil genannten Commits festgelegt.

Biologische Bezeichnungen werden bewusst als kompatible Profile
formuliert. Weder Antikörpermarker noch der auf der Karte sichtbare Ort
werden als alleiniger Beleg einer exakten Zellidentität behandelt.
Der Detailbericht dokumentiert sowohl die Ergebnisse als auch den
Schlussfolgerungsspielraum des vorhandenen Arbeitsstands.
\end{document}
